\documentclass[11pt,a4paper]{article}
\usepackage[T1]{fontenc}
\usepackage[utf8]{inputenc}
\usepackage{mathptmx}
\usepackage[margin=2.35cm]{geometry}
\usepackage{graphicx}
\usepackage{tikz}
\usetikzlibrary{arrows.meta,positioning,calc,backgrounds,fit}
\usepackage{amsmath}
\usepackage{booktabs}
\usepackage{microtype}
\usepackage{caption}
\usepackage{titlesec}
\usepackage[hidelinks]{hyperref}
\usepackage{xurl}
\usepackage{placeins}
\captionsetup{font=small,labelfont=bf}
\titleformat{\section}{\normalfont\large\bfseries}{\thesection.}{0.6em}{}
\titleformat{\subsection}{\normalfont\normalsize\bfseries}{\thesubsection}{0.6em}{}
\titlespacing{\section}{0pt}{1.1em}{0.4em}
\setlength{\parskip}{2pt}
\newenvironment{apabib}{\raggedright\begin{list}{}{%
  \setlength{\leftmargin}{1.6em}\setlength{\itemindent}{-1.6em}%
  \setlength{\itemsep}{3pt}\setlength{\parsep}{0pt}\setlength{\topsep}{4pt}}}%
  {\end{list}}
\newcommand{\Clm}{\mathrm{Cl}^{-}}
\newcommand{\ECl}{E_{\mathrm{Cl}}}

\title{\vspace{-1.6em}\textbf{From homeostasis to computation:\\[2pt]
could the GABA switch mature the core of a signed-XOR motif?}}

\author{Jes\'us Marco de Lucas\\[2pt]
\small Instituto de F\'isica de Cantabria (IFCA), CSIC--Universidad de Cantabria, Santander, Spain\\
\small \texttt{jesus.marco@csic.es} \quad ORCID:~0000-0001-7914-8494}
\date{\small \textsc{Hypothesis and Theory Article} \\ Preprint, August 2026}

\begin{document}
\maketitle
\vspace{-1.4em}

\begin{abstract}\noindent
We have proposed the \emph{signed-XOR mechanism}, a feedforward-inhibition motif in which a common inhibitory hub suppresses coincident drive to generate an XOR-type mismatch signal, while two opponent output channels convert that mismatch into a local \emph{signed} error, as a recurrent connectomic motif (Pe\~na Fern\'andez et al., 2026a), as a computational primitive (2026b), and as relevant for memory (2026c). In trying to understand why inhibitory neurons are so central in the hippocampus, a key site for memory, we were struck by the developmental ``GABA switch'': the same GABAergic conductance that is depolarizing and, under appropriate conditions, \emph{network-recruiting} in the early vertebrate brain becomes inhibitory and \emph{computation-controlling} upon maturation, when the chloride extruder KCC2 lowers intracellular chloride and moves the GABA reversal potential from above rest towards, and in many cells below, it. This note asks whether that switch could help to explain how a motif as intricate as the signed-XOR, an inhibitory hub that must sit at the centre of the excitatory arms it regulates, flanked by two opponent output channels, could \emph{emerge}. We develop a single, deliberately speculative thesis: \emph{computation is repurposed homeostasis}. 
We make the idea concrete in two conductance-based simulations: a single moving reversal potential flips a network from a synchronous builder to a stable governor, and, given a pre-positioned inhibitory hub, calcium-dependent
plasticity lifts the input-to-lateral arms from a subthreshold into a functional regime, after which the same hub, now inhibitory, expresses coincidence suppression and the read-out becomes XOR-like, with the supralinearly recruited inhibitory conductance acting at once as the mature operation and as the teaching signal. We are explicit about what this does and does not show: the architecture is preconfigured, so what the simulations demonstrate is the \emph{maturation} of a pre-existing scaffold rather than its construction de novo. The wider constructive claim is offered as a hypothesis, not as a result.
The chloride-controlled sign of GABA turns out to run deeper than vertebrate development: it is conserved down to \emph{C.~elegans}, though the worm shows the switch, not the computational motif, while the molecular kit of such a conductance-based governor is in fact pre-neuronal: plants already use ion gradients and gated anionic conductances for turgor and movement. What animals add is not chloride transport as such, since the cotransporter family is older than neurons, but an animal plasma-membrane chloride-\emph{extrusion} regime, later specialized in vertebrate neurons through KCC2, that moves the resting GABA reversal potential into an inhibitory range. Near rest this conductance acts predominantly as a shunt; at more negative reversal potentials it combines divisive and hyperpolarizing components. The sharpest test of the proposal is quantitative and local: the inhibitory conductance recruited onto a lateral by two coincident inputs must more than double relative to one, and by a margin large enough to satisfy the quantitative condition we derive.
\end{abstract}

\noindent\textbf{Keywords:} signed-XOR motif; chloride homeostasis; GABA switch; KCC2/NKCC1; conductance-based inhibition; shunting inhibition; opponent coding; BCM; predictive coding; evolution of computation.

\section{Introduction}\label{sec:intro}

We have argued in a series of recent contributions that a small, concrete circuit, the \emph{signed-XOR mechanism}, recurs in many connectomes and is a natural unit of neural computation: the idea began as a machine-learning reading of engram formation, in which XOR appears as a basic motif (Marco de Lucas et al., 2024), and as a homeostatic account of learning across species (Marco de Lucas, 2024); we then introduced the motif as an over-represented connectomic pattern for local directional error signalling (Pe\~na Fern\'andez et al., 2026a), framed its XOR operation as a neuronal loss function and computational primitive (Pe\~na Fern\'andez et al., 2026b), and related it to memory (Pe\~na Fern\'andez et al., 2026c). In this motif, two inputs each excite a ``lateral'' principal cell and, in parallel, a common inhibitory \emph{hub}; the hub returns conductance-based inhibition to these laterals and is recruited \emph{supralinearly} by coincident input, so that a read-out pooling the laterals responds strongly to a single input but is suppressed when both coincide, implementing an XOR-type gate, categorical once a downstream firing threshold is applied. Near a resting GABA reversal potential this inhibition has a clean shunting/divisive interpretation, but the conductance-based simulations below show that the gate survives, and becomes more selective, when the reversal is more hyperpolarizing and a subtractive component is added. Two aggregators of \emph{opposite neurotransmitter identity}, $F^{+}$ and $F^{-}$, then convert the mismatch into a local \emph{signed} error $e=F^{+}-F^{-}$, precisely the quantity a backpropagation-free learning rule needs. Thus the XOR-like coincidence suppressor and the mechanism that supplies its sign are distinct parts of the motif.  

This note takes a different perspective. In trying to understand \emph{why} inhibitory neurons should be so central in the hippocampus, a key site for memory, we arrived at a developmental fact that is well known but, in this context, striking: the \emph{GABA switch}. Early in the development of neural circuits, GABA does not act as an inhibitory neurotransmitter but as a \emph{depolarizing}, network-recruiting one (depolarizing, and under the right conditions excitatory). We suggest that the same interneurons that will later regulate the network may first help to \emph{build} it; only as KCC2, an electroneutral K$^{+}$--Cl$^{-}$ cotransporter that uses the outward potassium gradient to extrude chloride, rises in expression and transport activity does intracellular chloride fall, the GABA reversal potential drop towards rest, and the very same GABAergic conductance shift from a depolarizing builder to a mature inhibitory governor. Near rest that governor is predominantly shunting; below rest it acquires an additional hyperpolarizing component. One conductance; two opposite developmental roles, selected by a single homeostatic set point.

This duality is suggestive, because the signed-XOR motif is itself hard to specify from scratch: an inhibitory element that must occupy the centre of a convergent excitatory circuit, flanked by two output channels of opposite sign. Could the same switch that flips inhibition from builder to governor also help to explain how such a motif \emph{emerges}? The claim we develop is limited to the conductance-based \emph{coincidence-suppression core} of the motif: the pre-positioned inhibitory hub and its supralinear recruitment, which carve an XOR-type read-out. Shunting is the analytically clean near-rest limit of that inhibition, not the defining mechanism of the complete signed-XOR. The opponent pair $F^{+},F^{-}$ that turns the gate into a signed error is taken here as given, not explained. 

We develop a speculative thesis: \emph{computation is repurposed homeostasis}, as the acquisition of active chloride extrusion turns an ancient anionic turgor/movement signal, already present in plants\footnote{A clear caveat: the plant and animal molecules involved are \emph{not} homologous; this is functional repurposing of an ancient ion-control principle, not shared ancestry.}, into an inhibitory feedback governor, and the signed-XOR mechanism is one circuit in which that governor could be housed.  

To provide a visual scheme of the idea to the interested reader, the mature topology and operations of the motif are shown in Fig.~\ref{fig:mature}, while the two developmental steps that mature it thanks to the GABA switch are sketched in Fig.~\ref{fig:devseq}.


\definecolor{excgreen}{HTML}{2E7D32}
\definecolor{inhred}{HTML}{A11A1A}
\definecolor{latblue}{HTML}{1F4E79}
\definecolor{inpgrey}{HTML}{555555}
\definecolor{outpurple}{HTML}{5A3A8A}
\definecolor{nodefill}{HTML}{F4F4F8}

\begin{figure}[t]
\centering
\begin{tikzpicture}[
  every node/.style={font=\small},
  exc/.style={-{Stealth[length=2.4mm]}, line width=0.9pt, draw=excgreen},
  drive/.style={-{Stealth[length=2.4mm]}, line width=0.9pt, draw=inpgrey},
  inh/.style={-{Circle[length=2.0mm,fill=inhred]}, line width=1.0pt, draw=inhred},
  err/.style={-{Stealth[length=2.0mm]}, line width=0.7pt, draw=black!65},
  nin/.style={circle, draw=black!55, fill=inpgrey!85, text=white,
              minimum size=7.5mm, inner sep=0pt, font=\small\bfseries},
  nlat/.style={circle, draw=black!55, fill=latblue, text=white,
               minimum size=8mm, inner sep=0pt, font=\small\bfseries},
  nhub/.style={circle, draw=black!55, fill=inhred, text=white,
               minimum size=8.5mm, inner sep=0pt, font=\scriptsize\bfseries},
  nout/.style={circle, draw=black!55, fill=outpurple, text=white,
               minimum size=8.5mm, inner sep=0pt, font=\scriptsize\bfseries},
  nfp/.style={circle, draw=black!55, fill=excgreen, text=white,
              minimum size=7.5mm, inner sep=0pt, font=\small\bfseries},
  nfm/.style={circle, draw=black!55, fill=inhred, text=white,
              minimum size=7.5mm, inner sep=0pt, font=\small\bfseries},
  zone/.style={rounded corners=3pt, draw=black!30, dashed, line width=0.6pt},
]

% ---- nodes -------------------------------------------------------------
\node[nin]  (I)   at (0,  2.05) {I};
\node[nin]  (P)   at (0, -2.05) {P};
\node[nlat] (Li)  at (2.75,  2.05) {$L_i$};
\node[nlat] (Lp)  at (2.75, -2.05) {$L_p$};
\node[nhub] (H)   at (2.75,  0)    {INH};
\node[nout] (X)   at (5.55,  0)    {XOR};
\node[nfp]  (Fp)  at (8.35,  2.05) {$F^{+}$};
\node[nfm]  (Fm)  at (8.35, -2.05) {$F^{-}$};

% ---- excitatory arms and hub recruitment -------------------------------
\draw[exc]   (I) -- (Li);
\draw[exc]   (P) -- (Lp);
\draw[drive] (I) -- (H);
\draw[drive] (P) -- (H);

% ---- the hub inhibits the laterals only, never the read-out ------------
\draw[inh] (H) -- (Li) node[midway, left=1.2mm, inhred, font=\small] {$\div$};
\draw[inh] (H) -- (Lp) node[midway, left=1.2mm, inhred, font=\small] {$\div$};

% ---- laterals pool into the read-out -----------------------------------
\draw[exc] (Li) -- (X);
\draw[exc] (Lp) -- (X);

% ---- opponent read-out: the part that supplies the sign ----------------
\draw[exc] (X)  -- (Fp);
\draw[exc] (X)  -- (Fm);
\draw[exc, bend left=16]  (Li) to (Fp);
\draw[exc, bend right=16] (Lp) to (Fm);

\node[draw=black!45, fill=nodefill, rounded corners=2pt, inner sep=3.5pt,
      font=\small] (E) at (8.35, 0) {$e=\rho(F^{+})-\rho(F^{-})$};
\draw[err] (Fp) -- (E);
\draw[err] (Fm) -- (E);

% ---- the two operations, named separately ------------------------------
\begin{scope}[on background layer]
  \node[zone, fit=(I)(P)(Li)(Lp)(H)(X), inner sep=6pt] (zA) {};
  \node[zone, fit=(Fp)(Fm)(E), inner sep=6pt] (zB) {};
\end{scope}
\node[anchor=south, font=\footnotesize\bfseries, text=black!70]
      at ($(zA.north) + (0,1.2mm)$) {coincidence suppression: \emph{mismatch}};
\node[anchor=south, font=\footnotesize\bfseries, text=black!70]
      at ($(zB.north) + (0,1.2mm)$) {opponent read-out: \emph{sign}};

% ---- operating cases ---------------------------------------------------
\node[anchor=west, align=left, font=\scriptsize, text=excgreen]
      at (10.35, 1.15) {\textbf{one input on}\\ hub weakly recruited,\\
                        lateral escapes\\ $\to$ read-out ON};
\node[anchor=west, align=left, font=\scriptsize, text=inhred]
      at (10.35, -1.15) {\textbf{both inputs on}\\ hub recruited\\
                         \emph{supralinearly}\\ $\to$ read-out OFF};

% ---- legend ------------------------------------------------------------
\node[anchor=north west, align=left, font=\scriptsize] at (-0.75, -3.35)
  {\textcolor{excgreen}{\rule[0.6mm]{4.5mm}{0.7pt}}~excitatory \quad
   \textcolor{inpgrey}{\rule[0.6mm]{4.5mm}{0.7pt}}~excitatory drive onto hub \quad
   \textcolor{inhred}{\rule[0.6mm]{4.5mm}{0.9pt}$\bullet$}~conductance-based
   inhibition ($\div$: near-rest shunting limit)};
\end{tikzpicture}

\caption{\textbf{The mature signed-XOR motif and its operations.} Each input
excites its lateral and, in parallel, the common hub, which returns
conductance-based inhibition only to the two laterals. The figure draws the
analytically clean near-rest limit as a shunt ($\div$). A single input leaves
the hub weakly recruited, so its lateral largely escapes inhibition (read-out
ON); two inputs recruit the hub \emph{supralinearly}, so that inhibition grows
faster than the drive and suppresses both laterals (read-out OFF), a soft
XOR-type gate that a downstream firing threshold makes categorical. The
read-out projects to both aggregators, and each lateral to one of them
($L_i\!\to\!F^{+}$, $L_p\!\to\!F^{-}$); the two aggregators have opposite
neurotransmitter identity, and their difference is the local signed error
$e=\rho(F^{+})-\rho(F^{-})$. The two boxed operations are distinct and are the
reason the motif is more than a gate: the central inhibition generates mismatch
selectivity, and the opponent read-out supplies its direction.}
\label{fig:mature}
\end{figure}

\begin{figure}[t]
\centering
\includegraphics[width=\linewidth]{motif_dev_en.png}
\caption{\textbf{The developmental sequence: one inhibitory hub, three roles.} \textbf{(A) Builder.} While intracellular chloride is high, GABA is depolarizing ($+$) and the hub coactivates (recruits) the laterals $L_i,L_p$; GDP-like bursts grow the excitatory input arms (dashed). \textbf{(B) The switch.} As KCC2 rises and chloride falls \emph{in the lateral (postsynaptic) cells}, their GABA reversal potential $\ECl$ drops from above the firing threshold towards and, depending on cell and compartment, below rest, so that the hub's synapses onto them shift from depolarizing ($+$) to inhibitory ($-$). \textbf{(C) Governor.} The same hub, now central, suppresses the laterals; the schematic marks the near-rest shunting limit as $\div$, while more hyperpolarizing reversals add a subtractive component. The mature circuit and operations are detailed in Fig.~\ref{fig:mature}. Excitatory connections end in arrowheads; inhibitory ones, in filled circles.}
\label{fig:devseq}
\end{figure}

The following sections state the argument compactly; the results of the simulations and the supporting biology can be found in Appendices~\ref{app:sim}--\ref{app:plants}.

\section{The GABA switch and the motif it could implement}\label{sec:switch}
The GABA switch is a single change in the neurons that receive GABA: GABA$_{\mathrm{A}}$ receptors open a chloride channel, so that its reversal potential $\ECl$ is set by intracellular chloride in the postsynaptic cell, which the cotransporters actively control (its physiological value, $E_{\mathrm{GABA}}$, sits slightly above $\ECl$ because the channel also passes some bicarbonate; Appendix~\ref{app:bg}); NKCC1 imports chloride ($\ECl$ high, GABA depolarizing), whereas KCC2 extrudes it and lowers the GABA reversal potential into a mature inhibitory range. Neurons express NKCC1 early and KCC2 late, so that GABA begins depolarizing and becomes inhibitory, near-rest shunting in some cells and more hyperpolarizing in others, as has been observed in the hippocampus (Appendix~\ref{app:hipp}), but also in \emph{C.~elegans} (Appendix~\ref{app:worm}). This set point is not merely developmental drift: it is under active homeostatic control. The WNK--SPAK/OSR1 kinase cascade (a set of kinases that switch transporters on or off by phosphorylating them) reciprocally regulates the two cotransporters, activating NKCC1 while inhibiting KCC2, and WNK itself senses intracellular chloride, so that the very variable that decides whether GABA recruits or inhibits sits inside a molecular feedback loop (Appendix~\ref{app:bg}). The GABAergic \emph{conductance} is present the whole time; what changes is its driving force, and therefore whether opening that conductance depolarizes, predominantly shunts, or hyperpolarizes its target.

One premise must be flagged at the outset, because the whole argument rests on it. That immature GABA is \emph{depolarizing} is not disputed; that the depolarization is \emph{functionally excitatory in vivo}, and therefore able to build circuits, is. Slice results supporting a genuinely excitatory phase have been challenged as artefacts of the energy substrate or of the recording condition (Rheims et al., 2009; Bregestovski \& Bernard, 2012; Zilberter, 2016), and two-photon recordings in the intact neonatal neocortex find GABA depolarizing yet \emph{inhibitory} at the network level (Kirmse et al., 2015). Against this, spontaneous network events survive controls for mitochondrial energy metabolism (Ruusuvuori et al., 2010); optogenetic activation of GABAergic cells excites the neonatal network in vivo (Valeeva et al., 2016; Murata \& Colonnese, 2020); the perinatal oxytocin-triggered switch is itself an in vivo phenomenon (Tyzio et al., 2006); and interfering with depolarizing GABA in vivo impairs morphological and synaptic maturation (Cancedda et al., 2007; Wang \& Kriegstein, 2008). The dispute and its rebuttal are set out in Ben-Ari (2014) and the transporter biology is reviewed in Kaila et al. (2014). We adopt the constructive reading throughout, and we flag it as an \emph{assumption} rather than a datum: if the depolarizing phase proves non-constructive in vivo, the builder half of the present hypothesis fails with it, while the governor half, which concerns the mature inhibitory conductance, does not.

In its depolarizing phase an interneuronal hub \emph{coactivates} its targets and paces the synchronous bursts known as giant depolarizing potentials (GDPs), whose calcium transients grow the excitatory arms: the hub literally helps to build the circuit it will later regulate, and GABAergic \emph{hub} interneurons are known to orchestrate exactly this early synchrony (Bonifazi et al., 2009). After the GABA switch, the same hub, now central, \emph{inhibits} those targets: near rest it acts predominantly as a divisive shunt, whereas at more negative reversal potentials a hyperpolarizing/subtractive component is added. We call the two developmental modes \emph{builder} and \emph{governor} (Fig.~\ref{fig:devseq}, Appendix~\ref{app:hipp}).

The hardest question is topological: why should the inhibitory hub come to sit at the centre of the excitatory arms, wired so as to \emph{suppress} coincidence? Correlation-based growth cannot do it: Hebb's slogan, neurons which fire together wire together (Shatz, 1992), reinforces coincidence, which is exactly the opposite of what an XOR-type gate requires. But the switch supplies the missing \emph{anti-correlational} teaching signal, and the term matters: what is broken is not a symmetry between the two arms, which grow alike, but the equivalence between the single-input and the coincident condition. Development hands us a pre-positioned, broadly connected inhibitory hub (Fig.~\ref{fig:devseq}A); once that hub becomes inhibitory, ordinary calcium-dependent plasticity tunes the arms it already has into a functional regime. A lateral driven by a \emph{single} input sees only weak hub recruitment, depolarizes, admits a large calcium transient and \emph{potentiates} its arm; a lateral driven during \emph{coincidence} is clamped by the much larger inhibitory conductance, its calcium collapses, and its arm \emph{depresses}. The same conductance is therefore at once the mature coincidence-suppressing operation \emph{and} the teaching signal that tunes it; whether its instantaneous effect is mainly divisive or partly subtractive depends on the mature reversal potential.

Two conditions must hold for this to converge, and we state both as predictions. First, the inhibitory conductance delivered to a lateral must be \emph{supralinear} in the number of coincident inputs: writing $g_{\mathrm{inh}}(n)$ for the conductance recruited when $n$ inputs are active, two coincident inputs must draw more than twice the inhibition of one, $g_{\mathrm{inh}}(2)>2\,g_{\mathrm{inh}}(1)$, so that a single input recruits it weakly while a coincident pair clamps the lateral hard. Supralinearity in this weak form is necessary but not sufficient; the quantitative condition the read-out actually requires is Eq.~\eqref{eq:xorcond} below, which demands that the excess be large compared with the semi-saturation constant. In the present proof-of-principle model its source is the hub's firing threshold: coincident inputs recruit disproportionately more interneurons than a single input does. Whether a graded, non-spiking element could supply the same supralinearity is left open. A second non-linearity acts on the \emph{plasticity} rather than on the gate, and the two should not be conflated: the voltage dependence of dendritic calcium entry, through the Mg$^{2+}$ block of NMDA receptors and the sigmoidal activation of voltage-dependent calcium channels, makes the modest, inhibition-induced drop in depolarization collapse Ca$^{2+}$ influx steeply, carrying coincidence from potentiation into depression. In the proof-of-principle model this voltage-dependent calcium non-linearity is represented abstractly, by spike-triggered calcium increments; explicit NMDA-receptor and voltage-gated calcium-channel dynamics are not simulated, and \S\ref{app:switchsim} sets out what that abstraction costs. The first condition is about conductance and decides whether the mature circuit \emph{computes}; the second is about calcium and decides whether the developing circuit \emph{learns}. Second, the arms must \emph{clear} the firing ``cliff'', where a single input drives the lateral, but this need not be assumed: the depolarizing \emph{builder} phase lifts them there (Appendix~\ref{app:sim}). That the arms should begin below the cliff is not an arbitrary modelling choice but the expected starting condition of an immature projection, whose synapses are weak or functionally silent until activity recruits AMPA receptors (Isaac et al., 1995; Liao et al., 1995; Durand et al., 1996). Conductance-based simulations confirm that the full builder$\to$switch$\to$governor sequence is required, and that deleting any one element abolishes the result: with mature inhibition present from the start, with no hub at all, or with plasticity frozen, the arms never leave their initial value (\emph{failure to teach}); with a builder that never switches off, the circuit computes OR. The reason is not the one we first assumed: a control reported in Appendix~\ref{app:growth} shows that the arms it grows still compute XOR when probed at mature $\ECl$, so what that condition lacks is the \emph{sign} of the conductance rather than restraint in the size of the arms. What the sequence supplies is therefore not the wiring diagram, which the model is given, but the operating point: under a pre-positioned inhibitory governor, the developmental sequence moves the preconfigured circuit into a functional weight regime through ordinary calcium-dependent plasticity.

\section{From coincidence suppression to a signed-XOR computational motif}\label{sec:compute}
What development leaves behind is a topology, not yet a computation. The first computation is an XOR-like \emph{mismatch} signal, and what converts the topology into that signal is \emph{supralinearly recruited conductance-based inhibition}. Shunting is the analytically clean near-rest limit of this mechanism, not its defining property. It is useful to keep four steps distinct. \emph{(i)}~A common inhibitory hub returns a conductance $g_{\mathrm{inh}}$ to both laterals (\S\ref{app:bg}). \emph{(ii)}~When the inhibitory reversal potential lies near rest, that conductance acts predominantly as a shunt and admits the familiar divisive idealisation: the excitatory drive $D_i$ to a lateral is divided by a denominator that grows with hub recruitment. Writing $g$ for the hub's input--output recruitment function, the read-out of lateral $i$ is
\begin{equation}
y_i \;=\; \frac{D_i}{\sigma \;+\; \beta\, g\!\left(\textstyle\sum_{k} D_k\right)},
\label{eq:norm}
\end{equation}
the common inhibitory conductance appearing in the denominator. \emph{(iii)}~Coincidence is suppressed only if inhibition grows \emph{faster} than the drive, because a \emph{linear} pool ($g(x)=x$) does not suffice: with one input, $y_{10}=D/(\sigma+\beta D)$; with two, $y_{11}=2D/(\sigma+2\beta D)$, and $y_{11}/y_{10}=2(\sigma+\beta D)/(\sigma+2\beta D)>1$ for any $\sigma>0$. Linear divisive normalization can saturate the pooled read-out but cannot invert it: it yields gain control, not XOR. Suppression of an \emph{individual} lateral under coincidence is automatic for any increasing $g$; the non-trivial requirement is that this per-lateral suppression survive the pooling of two laterals rather than one, and therefore beat a factor of two. In the purely divisive limit, coincidence suppression requires $g(2D)>2\,g(D)$ and specifically
\begin{equation}
\beta\,\big[\,g(2D) - 2\,g(D)\,\big] \;>\; \sigma,
\label{eq:xorcond}
\end{equation}
so that doubling the input more than doubles the effective inhibitory denominator and drives the pooled read-out down. \emph{(iv)}~This graded coincidence suppression becomes an \emph{XOR-type} operation only after thresholding: a firing threshold $z=H(y-\vartheta)$ ($H$ the Heaviside step) turns the preference for exactly one active input into a categorical output.

The divisive step deserves an explicit caveat, since it is where the terminology can otherwise become misleading. That a near-rest inhibitory conductance divides a \emph{subthreshold} response is uncontroversial; that it divides the \emph{firing rate} is not. In a deterministic point neuron with a spike threshold, shunting inhibition acts largely subtractively on the $f$--$I$ curve rather than divisively (Holt \& Koch, 1997), and division emerges only when further ingredients are present: background synaptic noise, which smooths the threshold non-linearity (Chance et al., 2002; Prescott \& De Koninck, 2003), or dendritic saturation and the interaction of excitatory and inhibitory conductances in a spatially extended cell (reviewed in Silver, 2010). Carandini \& Heeger (2012) themselves record that the original shunting account of normalization did not survive contact with the data. Our simulations sit in the noisy, conductance-driven, near-threshold regime in which a divisive reading is defensible, but the chloride sweep of \S\ref{app:switchsim} shows that the computation does not depend on that special limit. The \emph{analytical idealisation} is a purely divisive shunt, Eqs.~\eqref{eq:norm}--\eqref{eq:xorcond}; the conductance-based model actually exhibits \emph{supralinearly recruited inhibition} whose divisive and hyperpolarizing/subtractive components vary with $\ECl$. Both components are delivered by the same hub and are therefore preferentially recruited under coincidence. Selectivity in fact improves as the mature reversal becomes more hyperpolarizing. The mechanism should therefore be named for what is common across the working range: \emph{supralinear conductance-based inhibitory coincidence suppression}.

We claim only the chain
\[
\text{supralinear inhibitory recruitment}
\;\longrightarrow\;
\text{coincidence suppression}
\;\longrightarrow\;
\text{thresholded XOR},
\]
and not a Boolean XOR at the biophysical level. The gate is graded (in our simulations, single-input drive exceeds coincident drive several-fold, not infinitely), so that ``exactly-one'' always refers to the thresholded read-out. The supralinearity is not an extra assumption bolted on for this section; it is the same condition the growth model needs to converge (\S\ref{sec:switch}), and in the circuit it comes from the hub's own firing threshold and population recruitment.

The second computation is the \emph{sign}, and it should not be conflated with the inhibitory mechanism that creates the XOR-like mismatch. A thresholded XOR reports only \emph{that} the two inputs disagree; it is a mismatch magnitude without a direction. Two aggregators of \emph{opposite neurotransmitter identity}, $F^{+}$ and $F^{-}$, read the circuit, and their difference
\begin{equation}
e \;=\; \rho(F^{+}) - \rho(F^{-}),
\label{eq:signederror}
\end{equation}
with $\rho$ a non-negative activation, is a local \emph{signed} error. The read-out projects to both aggregators while each lateral projects selectively to one of them ($L_i\!\to\!F^{+}$, $L_p\!\to\!F^{-}$), so that the residual difference carries not just mismatch but polarity. Because $\rho$ is non-linear, the common read-out need not cancel identically in the difference; it can instead gate the lateral-driven asymmetry and let the error vanish when the XOR-like read-out is silent. That is precisely the quantity a backpropagation-free learning rule needs, and the one that schemes such as equilibrium propagation (Scellier \& Bengio, 2017) and predictive coding (Whittington \& Bogacz, 2017) construct by other means; it is also the role we have argued the motif plays in local credit assignment (Pe\~na Fern\'andez et al., 2026a) and as a neuronal loss function (2026b).

The distinction has a direct biophysical meaning. Suppose, in the most local implementation, that the opponent outputs arrive as two conductances $g_{+}$ and $g_{-}$ onto the same membrane compartment, with reversal potentials on opposite sides of rest, $E_{+}=E_{\mathrm{L}}+\Delta E$ and $E_{-}=E_{\mathrm{L}}-\Delta E$. Linearising the membrane equation about rest gives
\begin{equation}
C_m\frac{dV}{dt} \;\approx\; \Delta E\,\big(g_{+}-g_{-}\big),
\label{eq:pushpull}
\end{equation}
so the direction of the charge or voltage excursion reports which opponent pathway dominates. A \emph{pure shunt} cannot provide this sign: if the inhibitory reversal lies at rest, opening that branch can attenuate a depolarisation but cannot drive a negative excursion away from rest. Thus the central hub and the opponent read-out perform different operations. The hub establishes \emph{whether} there is a mismatch by suppressing coincidence; the opponent pair establishes \emph{which way} the mismatch points. Equation~\eqref{eq:pushpull} is a possible membrane-level realisation of the latter, not a reduction of the complete multi-neuron motif to one compartment, and it imposes a stronger biophysical requirement than the abstract difference in Eq.~\eqref{eq:signederror}: a literal local signed voltage requires opponent driving forces, not shunting alone.

Nor need the signed error be transported far. In a local feedback reading, $F^{+}$ and $F^{-}$ could terminate on the same postsynaptic compartment, where the resulting push--pull voltage or calcium excursion sets the sign of plasticity (Appendix~\ref{app:sim}). The present simulations do not model this opponent membrane implementation explicitly; they establish the coincidence-suppression core. The complete signed-XOR should therefore be understood as two coupled operations---a conductance-based XOR-like mismatch gate and an opponent signed read-out---rather than as a single ``shunting mechanism''.

\section{An evolutionary reading}\label{sec:evo}

The materials used are ancient and widespread: plants also use chloride and anion channels to drive turgor and movement, and even use GABA to modulate anion channels, but they keep chloride \emph{high} and never acquire an animal-style plasma-membrane extrusion regime, so that they have the ancient ingredients without the specialization that places the GABA reversal potential in a mature inhibitory range (Appendix~\ref{app:plants}). On this reading, the GABA switch is one animal route by which an anionic conductance can be repurposed from drive into control.

Read broadly, the switch is an instance of a general move. \emph{Homeostasis} keeps the conditions for life; \emph{control} is homeostasis with feedback; and a particularly economical feedback element in an excitable membrane is a regulated conductance whose reversal potential lies at or below rest. Near rest it lowers input resistance and acts mainly as a divisive shunt; farther below rest the same opening also carries a hyperpolarizing current. Acquiring active chloride extrusion is what moves an anionic signal from a depolarizing regime into this inhibitory control range. The extrusion machinery (the cation-chloride cotransporter family) long predates neurons (Hartmann et al., 2014), so the innovation cannot be chloride transport as such; what is added is an animal \emph{plasma-membrane} chloride-extrusion regime, deployed in muscle as well as in neurons and later specialized in vertebrate neurons through KCC2, that holds the resting anionic driving force near rest. Nervous systems plausibly arose to coordinate contractile and ciliary effectors (Keijzer et al., 2013; J\'ekely et al., 2015), and their phylogenetic origin may not even be single (Moroz et al., 2014), so that the appearance of the chloride switch even in the \emph{muscle} of \emph{C.~elegans} fits a picture of conductance-based inhibitory control as a general tool of the excitable cell, deployed in muscle for movement and in neurons for computation. Plants provide an informative functional contrast: they use chloride, anion channels, GABA-related signalling, and turgor-based excitability, but they do not deploy an animal-style plasma-membrane chloride-extrusion regime to move GABA into the same mature inhibitory operating range. The parallel is convergent, not homologous: the conserved threads are the ancient cotransporter family and the \emph{strategy} of using an anionic conductance to set excitability, not a shared receptor.

\section{Predictions and caveats}\label{sec:pred}
\textbf{Predictions.} The sharpest remains quantitative, local, and available with existing methods. \emph{(1) Supralinear recruitment.} Voltage-clamping a lateral principal cell while stimulating one versus two convergent input pathways should yield an inhibitory conductance that more than doubles, $g_{\mathrm{inh}}(2)>2\,g_{\mathrm{inh}}(1)$, and in the near-rest divisive limit should satisfy Eq.~\eqref{eq:xorcond}; a linear or sublinear recruitment falsifies the coincidence-suppression core outright. In our own circuit the condition is met with a wide margin ($g_{\mathrm{inh}}(2)/g_{\mathrm{inh}}(1)=5.9$; Appendix~\ref{app:growth}), which fixes the quantity an experiment should be looking for and makes the prediction sharp rather than merely directional. \emph{(2) The pure shunt should not be a privileged point.} With the circuit and hub recruitment held fixed, shifting $E_{\mathrm{GABA}}$ from near rest into a moderately hyperpolarizing range should preserve coincidence suppression rather than abolish it; in the present model selectivity strengthens over the range tested because the newly added hyperpolarizing current is itself recruited supralinearly. A gate that works only at $E_{\mathrm{GABA}}\simeq E_{\mathrm{L}}$ would support the divisive idealisation but falsify the broader conductance-based account. \emph{(3) Shared conductance, shared fate.} Because the stabilizing and the XOR-computing functions are carried by the same GABAergic conductance, degrading chloride extrusion should degrade both together. Testing this cleanly requires care, because KCC2 also has ion-transport-independent structural roles in spine morphogenesis, so that a knockdown confounds the two; the experiment should therefore use a transport-dead KCC2 mutant, or an acute and reversible manipulation of the chloride gradient, with the structural function preserved as a control (Kaila et al., 2014). \emph{(4) Hub properties.} The inhibitory hub of a signed-XOR circuit should be early-maturing, broadly connected, and, less obviously, should receive substantially \emph{shared} afferents across hub cells: our sweeps show the gate failing as those afferents decorrelate, even while the arms still grow. \emph{(5) Calcium as teacher.} Blocking coincidence-induced postsynaptic calcium suppression during development should push the emergent read-out from XOR towards OR. \emph{(6) Compartmental grading.} More speculatively, if WNK-mediated regulation of KCC2 is local to the dendrite rather than somatic, the balance between shunting and hyperpolarizing inhibition should grade compartment by compartment, letting focal inputs on different branches sit at different chloride set points. \emph{(7) Opponent sign.} If $F^{+}$ and $F^{-}$ are realised locally as conductances onto one postsynaptic compartment, a literal signed voltage or calcium signal requires opponent driving forces, with the two effective reversal potentials on opposite sides of the local resting potential. Making the negative branch a pure near-rest shunt should preserve attenuation but abolish the negative excursion. This prediction belongs to the proposed membrane implementation of the signed read-out, not to the coincidence-suppression core established by the present simulations. \emph{(8) Duration and coupling should trade off.} The builder's requirement is a threshold in integrated drive rather than a minimum duration: in our sweeps, combinations of builder duration and hub-to-lateral coupling with the same product grow the arms to the same weight, to within $0.02$, up to about $450$~nS$\cdot$s (Appendix~\ref{app:growth}). Experimentally, shortening the depolarizing phase and strengthening the hub's coupling should therefore compensate one another, and a circuit deprived of half its depolarizing window should be rescuable by a proportionate increase in hub drive. We single this out because it does not follow trivially from the hypothesis and would be hard to accommodate if the builder acted by any mechanism other than accumulated calcium. A final expectation, that strong recurrent excitation \emph{before} the switch should give runaway epileptiform dynamics rather than computation, we list as a consistency check rather than a prediction, since it is already established that NKCC1-dependent depolarizing GABA facilitates seizures in the immature brain (Dzhala et al., 2005).

\textbf{Alternatives.} Coincidence suppression does not require this circuit, and the hypothesis should be judged against its competitors. Single neurons can compute linearly non-separable functions from passive dendritic properties alone (Caz\'e et al., 2013), and human layer 2/3 pyramidal dendrites implement an XOR-like operation intrinsically, through a graded, self-limiting dendritic calcium spike, with no inhibitory hub required (Gidon et al., 2020); interneuron dendrites are themselves non-linear compartments (Tzilivaki et al., 2019). Short-term synaptic depression at the input can likewise attenuate the coincident condition selectively. Feedforward inhibition, by contrast, is classically described as \emph{narrowing} the window for coincidence detection rather than suppressing coincidence (Pouille \& Scanziani, 2001), and it is worth being explicit that the signed-XOR hub is proposed to do the opposite. Inhibitory synaptic plasticity, finally, offers an independent route to a balanced, decorrelating operating point (Vogels et al., 2011). What distinguishes the present account is not that it is the only way to obtain an XOR-type read-out, but that it derives the read-out and its \emph{teaching signal} from a single conductance and ties both to a developmental variable that is independently measurable.

\textbf{Caveats.} The biology of each layer is established; the \emph{unifying} claim (a single governor repurposed from homeostasis to computation, and specifically the signed-XOR motif) is an interpretive hypothesis, not a result. The hypothesis addresses, strictly, the emergence of the conductance-based \emph{coincidence-suppression core}: the pre-positioned inhibitory hub and its supralinear recruitment, which carve an XOR-type read-out. The simulations do not establish the membrane-level implementation of the opponent signed read-out. Even there the simulations are weaker than the prose might suggest, and we prefer to say so plainly. The architecture is handed to the model, and what the developmental sequence contributes is the passage from a non-functional to a functional operating point: this is \emph{maturation}, not \emph{construction}. Nothing in our results explains why the hub should come to lie at the centre of the arms in the first place, which remains an assumption inherited from the connectomic observation rather than a result. And the constructive role of depolarizing GABA in vivo is contested (\S\ref{sec:switch}), so that the builder half of the argument stands or falls with it. It does \emph{not} explain the origin of the opponent read-out channels $F^{+}$ and $F^{-}$: why they should appear, why with opposite neurotransmitter identity, or why each lateral should project selectively to one of them. Those are taken here as part of the mature signed-XOR architecture proposed previously (Pe\~na Fern\'andez et al., 2026a,b); how the opponent pair itself develops remains an additional and unsolved problem. The plant/animal parallel is convergent. The worm's switch is largely cell-intrinsic, so that the activity-dependent, self-organizing version is a vertebrate elaboration, not a universal. The XOR is soft, not exact. And the growth model is a proof of principle with the limits set out in Appendix~\ref{app:sim}. We have tried to keep the established visibly separate from the speculative.

\section{Conclusion}\label{sec:concl}
The ingredients of a computational feedback element (chloride, anion channels, GABA, and the cotransporters that set the chloride gradient) are older than neurons and are already, in plants, part of the machinery of turgor and movement. What animals add is not chloride transport, which is older still, but an animal plasma-membrane chloride-extrusion regime, later specialized in vertebrate neurons through KCC2, that moves the anionic driving force from a depolarizing into an inhibitory range. Near rest the resulting conductance is predominantly shunting; at more negative reversal potentials it combines divisive and hyperpolarizing suppression. On this reading, the developmental GABA switch is not a curiosity of hippocampal wiring, but a route by which a hub is repurposed from builder to inhibitory governor, and, we suggest, by which a pre-existing signed-XOR scaffold could acquire a functional coincidence-suppression core. The full signed-XOR then requires a second operation, an opponent read-out that supplies direction to the mismatch; that sign cannot be attributed to shunting itself. We offer the synthesis as a target: to sharpen it, test it, and, where needed, knock it down.

\vspace{0.4em}
\noindent\textbf{Data and code.} The simulations (the GABA switch and the BCM-guided signed-XOR growth), the figure scripts, the random seeds used, and an environment specification pinning the Python and Brian2 versions are released at \url{https://github.com/jesusmarcodelucas/EmergingXOR} and accompany this preprint.

\vspace{0.2em}
\noindent\textbf{Competing interests.} The author declares a competing interest: a related European patent application, \emph{Signed XOR Feedback} (No.~26382252.0), has been filed (Marco de Lucas et al., 2026).

% ============================== APPENDICES ==============================
\FloatBarrier
\appendix
\titleformat{\section}{\normalfont\large\bfseries}{Appendix~\thesection.}{0.6em}{}
\titlespacing{\section}{0pt}{1.2em}{0.4em}


\section{Background biophysics, minimal simulations, and code}\label{app:sim}
\subsection{Background: driving forces, inhibitory conductances, and the calcium code}\label{app:bg}

\textbf{(i) Driving force.} A neuron's membrane holds a voltage $V$; each ion has an equilibrium (reversal) potential $E_{\mathrm{ion}}$, and the current through an open channel is the conductance $g$ times the driving force $(V-E_{\mathrm{ion}})$, that is, $I = g\,(V - E_{\mathrm{ion}})$. 
Opening a channel pushes $V$ towards $E_{\mathrm{ion}}$ \emph{and} lowers the input resistance. These are two inseparable consequences of adding conductance, but their relative importance depends on the reversal potential. When $E_{\mathrm{ion}}$ lies near rest, the channel changes $V$ little while lowering input resistance, so that it acts predominantly as a \emph{shunt} and can divide other inputs; when $E_{\mathrm{ion}}$ lies well below rest, the same conductance also carries a hyperpolarizing current and acquires a subtractive component. Thus ``shunting'' is an operating regime of conductance-based inhibition, not a different synaptic mechanism. The same channel can be excitatory, predominantly shunting, or hyperpolarizing depending on where $E_{\mathrm{ion}}$ sits relative to $V$.

\textbf{(ii) The chloride machinery.} GABA$_{\mathrm{A}}$ receptors open a chloride channel, so that their reversal potential is essentially $\ECl$, set by $[\Clm]_{\mathrm{in}}$. Concretely, because the open GABA$_{\mathrm{A}}$ channel is chloride-selective, $\ECl$ follows the Nernst relation
\begin{equation}
\ECl \;=\; \frac{RT}{z_{\Clm}F}\,\ln\frac{[\Clm]_{\mathrm{out}}}{[\Clm]_{\mathrm{in}}}
     \;=\; \frac{RT}{F}\,\ln\frac{[\Clm]_{\mathrm{in}}}{[\Clm]_{\mathrm{out}}},
\end{equation}
with valence $z_{\Clm}=-1$ for the anion; the sign is the point: opposite to a cation, \emph{raising} intracellular chloride moves $\ECl$ to more depolarized values (builder), while lowering it drags $\ECl$ towards rest (governor). The importer NKCC1 raises $[\Clm]_{\mathrm{in}}$ (pushing $\ECl$ above rest); the extruder KCC2 lowers it (dragging $\ECl$ towards rest). Immature neurons express NKCC1 early and KCC2 late, so that GABA depolarizes early and becomes inhibitory late; the mature effect is predominantly shunting when the reversal lies near rest and increasingly hyperpolarizing as it moves below rest. The \emph{GABA switch} is not a change of transmitter or receptor, but of a single homeostatic set point, $[\Clm]_{\mathrm{in}}$. (Strictly, the open GABA$_{\mathrm{A}}$ channel also passes bicarbonate, roughly $20\%$ as permeably as chloride, so the physiological reversal potential $E_{\mathrm{GABA}}$ follows the Goldman--Hodgkin--Katz equation and sits a few millivolts depolarized to $\ECl$ (Kaila, 1994); we write $\ECl$ throughout and collapse the two in the model. The family of transporters and its developmental regulation are reviewed by Kaila et al.\ (2014).)

Why these two proteins can hold such different set points is pure thermodynamics rather than biological detail. Both are electroneutral cotransporters (NKCC1 carries 1\,Na$^{+}$:1\,K$^{+}$:2\,$\Clm$, KCC2 carries 1\,K$^{+}$:1\,$\Clm$), so no net charge crosses the membrane and the free energy of transport contains no $zF\Delta\psi$ term. Direction is set by chemical potentials alone, and each transporter therefore has a hard equilibrium condition. KCC2 runs until $[\mathrm{K}^{+}]_{\mathrm{in}}[\Clm]_{\mathrm{in}}=[\mathrm{K}^{+}]_{\mathrm{out}}[\Clm]_{\mathrm{out}}$, that is, until $\ECl\rightarrow E_{\mathrm{K}}\approx-90$~mV. NKCC1 runs until $[\mathrm{Na}^{+}]_{\mathrm{in}}[\mathrm{K}^{+}]_{\mathrm{in}}[\Clm]_{\mathrm{in}}^{2}=[\mathrm{Na}^{+}]_{\mathrm{out}}[\mathrm{K}^{+}]_{\mathrm{out}}[\Clm]_{\mathrm{out}}^{2}$, that is, until $\ECl\rightarrow(E_{\mathrm{Na}}+E_{\mathrm{K}})/2\approx-15$~mV. The two transporters pull $\ECl$ towards attractors some 75~mV apart, and the neuron's actual $\ECl$ settles wherever their relative fluxes (together with the GABA$_{\mathrm{A}}$ and leak conductances) balance. The developmental switch is not a change of mechanism at all: it is a change in which attractor wins.

\textbf{(iii) Who moves the set point: the WNK--SPAK chloride thermostat.} The set point is not merely a passive consequence of gene expression; it is the output of a genuine homeostatic controller, and this is the level at which the ``computation is repurposed homeostasis'' thesis is most literal. The large cytoplasmic C-terminal domains of both cotransporters carry the phosphosites of the WNK--SPAK/OSR1 kinase cascade, and the cascade acts on them with \emph{opposite sign}: phosphorylation by SPAK/OSR1 activates NKCC1 (at Thr212/Thr217) and inhibits KCC2, while dephosphorylation does the reverse, so that a single kinase cascade reciprocally sets the balance of import and extrusion (Rinehart et al., 2009). Dephosphorylation is not merely the absence of kinase activity: the net transporter state reflects the balance between WNK--SPAK/OSR1 activity and opposing phosphatase activity, and maturation shifts that balance towards dephosphorylation of the inhibitory KCC2 sites. Pharmacologically shifting that balance dephosphorylates Thr$^{1007}$ of KCC2 together with the regulatory threonines of NKCC1 and the activation sites of SPAK itself, activating the extruder and inactivating the importer in concert (Zhang et al., 2020); no single dedicated phosphatase has been identified, and we leave the opposing activity unassigned. The two inhibitory threonines of KCC2, Thr906 and Thr1007, are highly phosphorylated in immature neurons and are progressively dephosphorylated as they mature; blocking or depleting WNK1 in immature neurons is by itself enough to enhance chloride extrusion and shift $E_{\mathrm{GABA}}$ towards hyperpolarizing values (Friedel et al., 2015). The loop closes because WNK is itself a \emph{chloride sensor}: $\Clm$ binds directly in the catalytic site and stabilizes the inactive conformation, inhibiting autophosphorylation (Piala et al., 2014). Low $[\Clm]_{\mathrm{in}}$ therefore activates WNK, which switches NKCC1 on and KCC2 off, and chloride rises; high $[\Clm]_{\mathrm{in}}$ inhibits WNK and the reverse occurs. Sensor, comparator and reciprocal actuators are all molecular, and together they implement a negative-feedback loop, a chloride thermostat, whose set point is exactly the variable that decides whether the inhibitory hub of the motif is a builder or a governor. On this reading the GABA switch is not a developmental accident overlaid on a computing circuit; it is a homeostatic controller being re-tuned, and the computation is what that re-tuning leaves behind.

\textbf{(iv) The calcium code.} Whether a synapse potentiates (LTP) or depresses (LTD) is governed, to good approximation, by postsynaptic calcium, which enters through voltage-dependent routes (NMDA receptors, voltage-dependent calcium channels) and therefore tracks how strongly the cell depolarized: \emph{moderate} calcium induces LTD, \emph{high} calcium LTP (Graupner \& Brunel, 2012; Clopath et al., 2010). The Bienenstock--Cooper--Munro (BCM) theory adds a \emph{sliding} threshold that follows the neuron's own recent activity, stabilizing the loop (Bienenstock et al., 1982); it is one member of a family of homeostatic mechanisms, alongside synaptic scaling, by which neurons hold their own excitability within a working range (Turrigiano, 2008). The consequence this note exploits is independent of whether mature inhibition is purely shunting: because calcium entry is voltage-dependent, a sufficiently recruited inhibitory conductance limits depolarization and controls the postsynaptic calcium of its targets, hence the sign of their plasticity. The same conductance that suppresses coincidence is therefore also a teaching signal.

\subsection{The switch as a moving reversal potential}\label{app:switchsim}
The claim that a single moving reversal potential turns a builder network into a governor is directly simulable. In a small recurrent conductance-based network (leaky integrate-and-fire neurons, with excitatory and GABAergic synapses and spike-frequency adaptation, in Brian2), we change \emph{only} $\ECl$ between conditions (Fig.~\ref{fig:sim}). With $\ECl$ depolarizing, recurrent GABA is excitatory and the network bursts into synchronous, GDP-like population discharges; with $\ECl$ in the mature inhibitory range, the identical circuit is sparse and stable; a slow $\ECl$ ramp reproduces the transition, with the discharges disappearing as the reversal potential crosses towards threshold. It is a proof of principle: the sign of a single homeostatic variable suffices to move a circuit between a cooperative, synchronizing regime and a stable, decorrelating one. It does not model synaptic growth; that is Appendix~\ref{app:growth}.

The chloride set points should be declared rather than left implicit. This network uses $\ECl=-45$~mV immature and $\ECl=-75$~mV mature, so that its mature condition is frankly \emph{hyperpolarizing}; the growth model of Appendix~\ref{app:growth} instead uses $-40$~mV and $-65$~mV, the latter sitting exactly at rest and therefore purely shunting. Taken together the two simulations bracket the mature regime, from hyperpolarizing to purely divisive, and the qualitative result (synchronous builder $\to$ sparse governor) holds at both ends, which is a modest robustness argument. What neither tests is the immature end: both place $\ECl$ \emph{above} the spike threshold, so GABA is frankly excitatory rather than merely depolarizing, whereas measured immature values are typically nearer $-50$ to $-45$~mV, that is, subthreshold, with excitation arising indirectly through relief of the Mg$^{2+}$ block and recruitment of voltage-dependent calcium channels. That choice is not innocuous, since it makes the builder maximally effective, and we have therefore mapped it.

The immature set point turns out to be bounded, and sharply. Sweeping $\ECl^{\mathrm{build}}$ from $-60$ to $-40$~mV at twenty seeds per value, with the mature value held at $-65$~mV, the builder fails completely at $-60$, $-55$ and $-50$~mV: in every one of sixty runs the arms remain at exactly $w_0=0.10$, the laterals never fire, postsynaptic calcium stays at zero, and the BCM rule of Eq.~\eqref{eq:bcm} returns $\Delta w=0$ identically. The failure is algebraic rather than statistical. While the arms sit at $w_0$ their excitatory conductance is negligible, so a lateral's steady-state potential is essentially a convex combination of rest and the chloride reversal,
\[
V \;\approx\; \frac{g_{\mathrm{L}}E_{\mathrm{L}} + g_{\mathrm{i}}\ECl}{g_{\mathrm{L}}+g_{\mathrm{i}}},
\]
which lies strictly between $E_{\mathrm{L}}$ and $\ECl$; if $\ECl \leq V_{T}$ the lateral cannot reach threshold for \emph{any} GABAergic conductance whatever, and no amount of hub recruitment will move it. With the measured $g_{\mathrm{i}}=44$~nS the boundary sits near $-48$~mV, and the sweep brackets it: $-50$~mV fails in $20/20$ runs, $-45$~mV succeeds in $20/20$.

The model therefore requires immature GABA to be \emph{frankly excitatory} and not merely depolarizing. We state this as a limitation rather than a detail, because it is exactly the point on which the in vivo literature divides (\S\ref{sec:switch}): if the depolarization of immature GABA is subthreshold in the intact brain, as Kirmse et al.\ (2015) report for neonatal neocortex, then the builder mechanism as modelled here does not run at all.

This boundary is, however, specific to the present spike-triggered calcium proxy, and we would not want it read as a general bound on the hypothesis. In the model, calcium increments only on a postsynaptic spike, so a lateral that never fires has exactly zero calcium and the BCM rule of Eq.~\eqref{eq:bcm} returns exactly zero weight change; the boundary at $\ECl=V_T$ then follows by construction. The biology we invoke elsewhere (\S\ref{app:bg}(iv)) offers at least three routes by which a \emph{subthreshold} depolarization could still admit calcium: relief of the Mg$^{2+}$ block of NMDA receptors, activation of voltage-dependent calcium channels below spike threshold, and summation with concurrent glutamatergic drive. None of these is implemented here. The sweep therefore bounds this implementation rather than the biological hypothesis, and it identifies the next simulation to run: a lateral with a continuously voltage-dependent calcium source, or a multicompartmental one, which would test whether physiologically subthreshold depolarizing GABA can supply the calcium that growth requires. We expect that it could, and that expectation is itself a prediction.

The sweep also returned something we did not anticipate. The published set point is not the best one. At $\ECl^{\mathrm{build}}=-45$~mV, closer to measured immature values than the $-40$~mV used elsewhere, every run grows both arms past the firing cliff (median $w=0.31$) and the read-out is \emph{more} selective than at $-40$~mV: $\mathrm{xor\_index}$ median $0.86$ (IQR $[0.82,1.00]$) against $0.78$. The circuit is correspondingly quieter, with single-input rates of ${\sim}21$~Hz against ${\sim}41$~Hz and coincidence held to $3$~Hz against $8$~Hz. The result is thus robust across the whole suprathreshold range tested, and the $-40$~mV we flagged above as a favourable choice turns out to be favourable for \emph{drive} rather than for selectivity.

The mature set point behaves quite differently, and it retires a concern rather than confirming one. Sweeping $\ECl^{\mathrm{gov}}$ from $-75$ to $-60$~mV with the builder held at $-40$~mV, the read-out computes XOR in $20/20$ runs at $-75$, $-70$ and $-65$~mV, and in $19/20$ at $-60$~mV, so the result survives across the entire range (Fig.~\ref{fig:eclsweep}B). We had flagged the published value, $\ECl^{\mathrm{gov}}=-65$~mV, as convenient because it sits exactly at rest and therefore makes inhibition purely divisive, which is the regime the algebra of \S\ref{sec:compute} assumes. The sweep shows the opposite of a convenient choice. Selectivity improves \emph{monotonically} as inhibition becomes more hyperpolarizing: the coincidence rate falls from $R_{11}=11.5$~Hz at $-60$~mV through $9.0$ at $-65$ and $7.0$ at $-70$ to $5.5$~Hz at $-75$, while single-input rates stay between $40$ and $44$~Hz throughout, so the gain is in suppression and not in a general silencing. The median $\mathrm{xor\_index}$ rises correspondingly, from $0.72$ to $0.82$.

The reason is worth stating, because it qualifies our own framing. A hyperpolarizing reversal adds a subtractive current on top of the divisive one, and that subtractive current is itself recruited supralinearly, since the conductance it flows through is about six times larger under coincidence than under a single input. Subtraction therefore acts preferentially where the gate needs it, and adds to the division rather than competing with it. Two consequences follow. The empirical result is more robust than the theory that motivated it, holding well outside the purely divisive regime. But the mechanism in the simulations is not purely divisive either, so the normalization algebra of Eqs.~(\ref{eq:norm})--(\ref{eq:xorcond}) should be read as an idealisation that isolates the divisive contribution, not as a complete account of what the conductance-based model does. We prefer to say this plainly: the pure-shunt set point is the middle of the working range, not its favourable end.

\begin{figure}[h]
\centering
\includegraphics[width=0.9\linewidth]{demo_brian2_en.png}
\caption{\textbf{The GABA switch as a moving reversal potential (Brian2).} The same excitatory/inhibitory microcircuit at three chloride set points. (A) Immature, depolarizing GABA: synchronous, GDP-like discharges. (B) Mature, inhibitory GABA: sparse, stable activity. (C) A developmental $\ECl$ ramp: the discharges disappear as the reversal potential crosses threshold. (D) The mean firing rate collapses across the switch. Only $\ECl$ differs between conditions.}
\label{fig:sim}
\end{figure}

\begin{figure}[h]
\centering
\includegraphics[width=\linewidth]{ecl_sweep_en.png}
\caption{\textbf{Chloride set-point robustness of the two-phase model} (full sequence, $n=20$ seeds per point; every run classified individually, and labelled on whether the arms cleared the firing cliff rather than on an absolute read-out rate). \textbf{(A) Builder.} Sweeping $\ECl$ during the depolarizing phase with the mature value held at $-65$~mV. Within the present spike-triggered-calcium implementation, below the spike threshold $V_T=-50$~mV the arms never move from $w_0=0.10$ in any of sixty runs: while the arms are subthreshold a lateral's potential is a convex combination of $E_{\mathrm{L}}$ and $\ECl$, so no GABAergic conductance can carry it to threshold. Above $V_T$ the builder teaches and the circuit computes XOR in $20/20$ runs at both $-45$ and $-40$~mV. \textbf{(B) Governor.} Sweeping $\ECl$ during the mature phase with the builder held at $-40$~mV. XOR survives across the whole range, and selectivity improves monotonically as inhibition becomes more hyperpolarizing: the coincidence rate $R_{11}$ falls from $11.5$ to $5.5$~Hz while single-input rates stay near $40$~Hz. The purely shunting set point used elsewhere ($-65$~mV, dotted) is the middle of the working range, not its favourable end.}
\label{fig:eclsweep}
\end{figure}

\FloatBarrier
\subsection{Growing the XOR-like read-out: the coincidence-suppression core}\label{app:growth}
Correlation-based growth cannot carve XOR: neurons that fire together wire together (Shatz, 1992), which reinforces coincidence, whereas XOR must \emph{suppress} it. The required anti-correlational teaching signal is supplied instead by the switch, through the calcium-dependent plasticity of \S\ref{app:bg}(iv). Because the hypothesis concerns the \emph{transition} from depolarizing to inhibitory GABA, $\ECl$ must \emph{move} during the run rather than be held fixed. We build the circuit in Brian2: two Poisson inputs, two lateral principal cells, a pre-positioned set of inhibitory hub neurons (fixed, broadly connected, projecting only to the laterals and not to the read-out, in keeping with Fig.~\ref{fig:mature}), and a read-out; only the direct input$\to$lateral arms are plastic. The two lateral cells should be read as the two principal-cell \emph{channels} of the motif (the left and right arms of the XOR), not as a realistic excitatory population size; a more biological version would replace each lateral by an excitatory assembly (for example ${\sim}40$ cells per arm) sharing a common inhibitory hub whose total conductance is normalized. Because identical copies of a lateral would behave identically, the value of such a population model lies in intra-arm \emph{heterogeneity} and partial assemblies rather than in cell number; we leave that control to future work.

The arms start \emph{below} the firing ``cliff'' ($w_0=0.10$, cliff $\approx0.21$), where the \emph{mean} single-input drive is subthreshold: a lateral spikes only rarely (at most a few hertz, from chance Poisson clusters), calcium is near zero, and plasticity is negligible on its own, so the arms cannot bootstrap unaided. This starting condition carries the necessity result and should therefore be justified rather than assumed: immature glutamatergic projections are characteristically weak or functionally silent, expressing NMDA receptors before AMPA receptors and becoming conducting only once activity recruits the latter (Isaac et al., 1995; Liao et al., 1995; Durand et al., 1996). Had the arms started above the cliff, the governor alone would suffice and the builder would be dispensable; the claim that the developmental sequence is necessary is thus conditional on a subthreshold start, which is what the biology of immature synapses supplies. The run has three phases: a \emph{builder} ($\ECl=-40$~mV, GABA depolarizing), a \emph{switch} ($\ECl$ ramps $-40\to-65$~mV over $4$~s), and a \emph{governor} ($\ECl=-65$~mV, here chosen at rest and therefore in the pure-shunt limit). The builder phase itself \emph{produces} the above-cliff starting condition rather than assuming it: while GABA is depolarizing the hub coactivates the laterals in GDP-like bursts that carry calcium and lift the arms over the cliff; the switch then converts the hub to inhibition; and the grown, now supra-threshold arms are used by the governor to compute XOR. When one input is active its lateral sees weak inhibition, gains high calcium and \emph{potentiates}; during coincidence the supralinearly larger inhibitory conductance clamps the lateral, calcium collapses, and the arm \emph{depresses}. In this particular run that conductance is purely shunting because $\ECl=E_{\mathrm{L}}$; the chloride sweep above shows that the same computational and teaching roles persist when a hyperpolarizing component is added.

Each plastic arm $w$ (input$\to$lateral) follows a BCM sliding-threshold rule, updated every $\Delta t=1$~ms:
\begin{equation}
\Delta w \;=\; \eta\,x\,c\,(c-\theta),\qquad w \leftarrow \mathrm{clip}(w,\,w_{\min},\,w_{\max}),
\label{eq:bcm}
\end{equation}
where $x$ is a presynaptic eligibility trace ($\dot{x}=-x/\tau_{\mathrm{pre}}$, $+1$ per presynaptic spike, $\tau_{\mathrm{pre}}=15$~ms); $c=\mathrm{Ca}/\kappa$ is a postsynaptic calcium proxy ($\dot{\mathrm{Ca}}=-\mathrm{Ca}/\tau_{\mathrm{Ca}}$, $+1$ per postsynaptic spike, $\tau_{\mathrm{Ca}}=60$~ms, $\kappa=9$); $\theta$ is the sliding threshold ($\dot{\theta}=(c^{2}-\theta)/\tau_{\theta}$, $\tau_{\theta}=800$~ms); $\eta=3\times10^{-5}$; and the weights are clipped to $[w_{\min},w_{\max}]=[0.05,2.6]$, with no explicit decay term (when $c=0$ the rule gives $\Delta w=0$). As in the original BCM formulation, the threshold tracks $\langle c^{2}\rangle$ while being compared against $c$; the comparison is consistent only because $c$ is a \emph{dimensionless} proxy whose scale is fixed by the normalizing constant $\kappa$, and $\theta$ should be read as a set point on that same normalized axis rather than as a quantity in units of $c^{2}$.

The supralinearity that Sec.~\ref{sec:compute} requires is a property of this architecture and can be read off it directly, so we measure it rather than assume it. Holding the arms fixed and probing the circuit with one versus two active inputs, the hub population rate rises from $6.8$ to $50.9$~Hz (a factor of $7.5$) and the inhibitory conductance delivered to a lateral rises from $g_{\mathrm{inh}}(1)=44$~nS to $g_{\mathrm{inh}}(2)=259$~nS, a factor of $5.9$. The weak condition $g_{\mathrm{inh}}(2)>2g_{\mathrm{inh}}(1)$ is therefore met with a wide margin, and so is the quantitative form, Eq.~\eqref{eq:xorcond}: the excess $g_{\mathrm{inh}}(2)-2g_{\mathrm{inh}}(1)=171$~nS greatly exceeds the effective semi-saturation term $\sigma_{\mathrm{eff}}=g_{\mathrm{L}}+g_{\mathrm{e}}\approx17$~nS. The reason is that the hub sits \emph{near its threshold} when a single input is active: with one input the mean hub drive is subthreshold and the cells fire only on fluctuations, while two inputs carry the population across threshold. This is precisely the recruitment non-linearity claimed in Sec.~\ref{sec:switch}, and it is worth stating that it is a property of the operating point rather than of the wiring: a hub driven hard enough to saturate with a single input recruits at best linearly, and the circuit then fails to suppress coincidence.

Two further details matter. First, the laterals carry a spike-frequency \emph{adaptation} current: without it, depolarizing GABA clamps a lateral at $\ECl=-40$~mV and it fires tonically at saturation, which drives $\theta=\langle c^{2}\rangle$ \emph{above} $c$, flips the sign of $(c-\theta)$, and makes the builder \emph{depress} the arms it should grow. Adaptation makes the depolarizing phase self-terminate into brief, high-calcium \emph{bursts} separated by quiet gaps, resembling the burst-and-gap temporal structure of giant depolarizing potentials, keeping $c>\theta$ so potentiation persists; the GDP structure is thus what lets the builder teach. Second, the test is hardened against artefacts: the hub cells carry \emph{heterogeneous} firing thresholds, so the bursts do not depend on artificial synchrony, and every probe follows a clean reset of all dynamic state, so residual adaptation cannot leak in.

\begin{figure}[h]
\centering
\includegraphics[width=0.86\linewidth]{phase_trace_en.png}
\caption{\textbf{The builder grows the arms; the governor uses them (Brian2, full sequence).} $\ECl$ (top) is depolarizing in the builder phase, ramps through the switch, and rests in the governor phase. The two direct input arms (bottom) start at $w_0=0.10$, \emph{below} the firing cliff (dashed): while GABA is depolarizing the hub coactivates the laterals in GDP-like bursts that carry calcium and lift the arms over the cliff; after the switch the grown, supra-threshold arms are used by the inhibitory governor to compute XOR (here at the pure-shunt set point). Without the builder the arms never leave $0.10$.}
\label{fig:trace}
\end{figure}

The developmental sequence is \emph{necessary}, not decorative, and a set of lesion controls shows it (Fig.~\ref{fig:factorial}); these are single-element deletions rather than a full factorial design. With mature inhibition present from the start (no builder), or with no hub at all, the laterals never fire and the arms stay at $0.10$ (\emph{failure to teach}). Running the \emph{same} $\ECl$ transition with plasticity frozen ($\eta=0$) collapses too, so the transition alone builds nothing: it is the plasticity, gated by the changing driving force of the same GABAergic conductance, that tunes the circuit. A builder that never switches off gives \emph{OR}, but the reason is not the one we first assumed, and the control is worth reporting because it corrects the mechanism rather than the result. That condition is probed with $\ECl$ still depolarizing, since by construction the switch never happens, so the hub is still \emph{adding} excitation under coincidence. When the arms it grew ($w\approx0.92$) are instead probed at mature $\ECl$ with the hub active, the circuit still computes XOR ($R_{10}=95$, $R_{11}=33$~Hz, $\mathrm{xor\_index}=0.62$), and does so with the same selectivity as the full sequence itself ($w\approx0.57$, $\mathrm{xor\_index}=0.62$). Over-growth alone does not produce OR at that weight; what produces it is that GABA never became inhibitory. Over-growth does eventually invert the read-out, but only well beyond the weights the builder reaches: probed at mature $\ECl$, $\mathrm{xor\_index}$ falls smoothly from $0.81$ at $w=0.45$ through $0.55$ at $w=1.20$ to $0.44$ at $w=1.80$, crossing our OR criterion between $1.2$ and $1.8$. The necessity of the switch is therefore established more strongly than before, and by a cleaner route: it is the inhibitory sign of the GABAergic current (set by its driving force), not the size of the arms, that the builder-only condition lacks. Only the full builder$\to$switch$\to$governor sequence gives \emph{XOR}, a single input driving the read-out several times harder than coincidence. A direct \emph{arm-acquisition} test (probing single inputs with the hub silenced) confirms the developmental claim: after the full sequence \emph{both} arms fire on their own, so the builder lifted them over threshold, whereas after any failure-to-teach condition \emph{neither} arm has grown. Equivalently, probed \emph{before} any training the circuit collapses, and only \emph{after} the full sequence does it compute XOR: the sequence, not either endpoint alone, brings the circuit into its functional regime.

\begin{figure}[h]
\centering
\includegraphics[width=\linewidth]{phase_factorial_en.png}
\caption{\textbf{Each developmental phase (and plasticity itself) is necessary (Brian2; arms start below the cliff, $w_0=0.1$; heterogeneous hub; clean reset before probing).} Read-out rates for the four input patterns ($10,01,11,00$) after training, under five manipulations. \emph{Governor only}, \emph{no hub}, and \emph{no plasticity} ($\eta=0$: the full $\ECl$ transition with learning frozen) $\to$ \emph{failure to teach} (arms stay at $0.1$). \emph{Builder only} (the switch never happens) $\to$ OR (coincidence, $11$, is largest). \emph{Full sequence} $\to$ XOR ($47,47,9,0$~Hz; arms grown from $0.10$ to $\approx0.57$). Beneath each panel: the trained arm weights and an \emph{arm-acquisition} verdict (single inputs probed with the hub silenced): ``both grown'' only where the builder ran, ``neither grown'' in every failure-to-teach case. Note on the \emph{builder only} panel: the OR verdict is obtained with $\ECl$ still depolarizing, since by construction the switch never happens. The same trained arms re-probed at mature $\ECl$ compute XOR ($\mathrm{xor\_index}=0.62$); what that condition lacks is the \emph{sign} of the conductance, not restraint in the growth of the arms.}
\label{fig:factorial}
\end{figure}

Finally, the \emph{timing} of the switch matters, but less simply than we first reported, and the route by which we found that out is worth stating. An earlier, exploratory sweep of $90$ runs (three seeds per cell, correlated hub) suggested a window bounded below by collapse and above by OR. Re-scoring it on the criteria set out below changed $26$ of its $90$ labels, and $22$ of those $26$ had been called \emph{collapse} while in fact computing XOR, with median $\mathrm{xor\_index}=0.86$; only four were authentic collapse, in the sense that the arms never left $w_0$. Its lower band was therefore very largely an artefact of the absolute rate gate rather than a developmental boundary. Rather than redraw a superseded sweep with better labels, we replaced it: Fig.~\ref{fig:window} reports a new run of the hardened, heterogeneous-hub model at twenty seeds per cell over an extended range of builder durations, $800$ runs in all.

The window that emerges is real but has two edges of quite different kinds.

The early edge is not a duration at all. Authentic failures to teach occupy a single corner of the grid, at the shortest builder and the weakest coupling, and what separates them from success is the \emph{product} of duration and coupling rather than either alone. Arms grown under different combinations with the same product reach the same weight to within $0.02$ for products up to about $450$~nS$\cdot$s: $(t_{\mathrm{build}},w_{\mathrm{hl}})=(25,16)$ and $(40,10)$ give $w=0.54$ and $0.52$, and $(15,16)$ and $(25,10)$ give $0.36$ and $0.34$. Beyond that the equivalence degrades to a mean discrepancy of $0.11$, as the sliding threshold saturates growth. The builder's requirement is therefore a threshold in \emph{integrated drive}, near $t_{\mathrm{build}}\times w_{\mathrm{hl}}\approx100$--$130$~nS$\cdot$s, and not a minimum duration. This yields a prediction we had not anticipated, and one that does not follow trivially from the hypothesis: shortening the depolarizing phase and strengthening the hub's coupling should compensate one another.

The two axes are not independent either, which we had not appreciated. The hub-to-lateral weight enters the \emph{same} conductance in both phases, since only the reversal potential changes between them, and while $\ECl$ is depolarizing that conductance is a drive rather than an inhibitor. Sweeping it therefore sweeps the strength of the builder at the same time as the strength of the mature inhibitory governor. The data show this directly: the arms, which grow only during the builder phase, scale by a factor of $1.9$ to $2.5$ across that axis at every builder duration. The abscissa of Fig.~\ref{fig:window} should accordingly be read as the hub-to-lateral coupling, not as mature inhibitory strength alone, and the apparent lower boundary is at least as much a builder-strength effect as a timing effect. We stress that this is a consequence of the hypothesis rather than a defect of the model: if one conductance really does serve both roles, then any parameter that scales it necessarily scales both, and the two axes of such a sweep cannot be independent. Separating the builder-phase from the governor-phase weight is therefore available to us only as a \emph{control}, one that asks which of the two roles produces which boundary, and not as a more correct version of the model.

The late edge, by contrast, \emph{is} a duration, and it is nearly independent of coupling. Selectivity declines monotonically with builder duration in every column, and crosses the OR criterion between $t_{\mathrm{build}}=50$ and $65$~s: at $50$~s the median $\mathrm{xor\_index}$ sits at $0.49$--$0.56$ across the five couplings and the cells are genuinely mixed, while at $65$~s four of the five columns are OR by majority and the fifth is an exact $10/10$ tie. Two cells of the sweep, at $(t_{\mathrm{build}},w_{\mathrm{HL}})=(50,16)$ and $(65,34)$, split their twenty seeds evenly between XOR and OR; we report them as ties rather than assigning them a majority, since at the boundary the categorical label is precisely what the data do not determine. That the crossing should fall at much the same duration whatever the coupling is intelligible, since coupling scales the mature inhibition as well as the arms that inhibition must control, so the two effects largely cancel; what does not cancel is the accumulated growth itself.

Two qualifications on that late edge. The transition is gradual, spanning some fifteen seconds, so it is a soft boundary, and the continuous index rather than the categorical label carries the information. And at the strongest coupling with the longest builder the arms approach the hard bound $w_{\max}=2.6$ (median $2.27$), so the anomalous survival of XOR in that single cell reflects the ceiling limiting further growth rather than a recovery of selectivity; we exclude it from the argument.

The developmental window is therefore real and has an interior optimum, as we originally claimed, but for different reasons on its two sides and over a broader plateau than the published sweep suggested: XOR is the majority outcome at every coupling from $t_{\mathrm{build}}=15$ to $32$~s, unanimous in most cells of that band and never below $18/20$ in any of them.

\begin{figure}[h]
\centering
\includegraphics[width=0.72\linewidth]{window2d_en.png}
\caption{\textbf{A developmental window for the switch} (Brian2, hardened heterogeneous-hub model, $20$ seeds per cell, $800$ runs). Outcome of the full two-phase sequence as the builder duration $t_{\mathrm{build}}$ is swept against the hub-to-lateral coupling $w_{\mathrm{HL}}$. Each cell gives the majority class, the vote, and the median arm weight and $\mathrm{xor\_index}$; labels gate on whether the arms cleared the firing cliff, never on an absolute read-out rate (\S\ref{app:growth}). The abscissa is the hub-to-lateral coupling and \emph{not} the mature inhibitory strength alone: the same conductance carries the depolarizing builder drive and the mature inhibition, so this axis scales both. Authentic failure to teach occupies one corner only, at the shortest builder and the weakest coupling; the late, OR edge lies between $t_{\mathrm{build}}=50$ and $65$~s and is nearly independent of coupling. The cell at the longest builder and strongest coupling is excluded from the argument, its arms having approached the hard bound $w_{\max}=2.6$.}
\label{fig:window}
\end{figure}

\textbf{Robustness across seeds, and the role of hub correlation.} Because a single run cannot support a developmental-learning claim, we repeated the full sequence over $100$ random seeds and classified \emph{each} run individually (Fig.~\ref{fig:robust}). With the shared-input hub the outcome is effectively deterministic: XOR in $100/100$ runs, both arms grown in $100/100$, median $\mathrm{xor\_index}=0.78$ (IQR $[0.70,0.81]$), coincidence held to $R_{11}=8$~Hz (IQR $[8,11]$) against single-input rates of ${\sim}41$~Hz; a fully-random (unbalanced) pattern-exposure control is essentially identical, so the result is not an artefact of balanced exposure. Four runs of that control were reported as exceptions because their single-input rates fell below the former absolute-rate gate at $25$~Hz; their continuous index was XOR-like throughout. The one manipulation that does matter is the correlation of the hub's afferents: when each hub cell receives a mixture of the shared trains and its own independent, pattern-driven train, coincidence suppression degrades smoothly as sharing is reduced: $R_{11}$ rises monotonically ($8\to87$~Hz), $\mathrm{xor\_index}$ falls and crosses zero near $\mathrm{shared}\approx0.66$, and the XOR fraction drops through a knee near $0.82$, becoming OR by $\mathrm{shared}\leq0.6$. Crucially, \emph{both arms grow in $100\%$ of runs across the entire sweep}: what fails as the hub de-correlates is not the growth of the direct arms but the coincidence-dependent inhibitory recruitment. This sharpens what the model shows: the plasticity does not discover the XOR architecture, which is preconfigured (both inputs onto the hub, the hub's supralinear threshold, the hub onto both laterals, the read-out pooling both laterals, mature GABA inhibitory); rather, the two-phase sequence \emph{matures} it, the depolarizing builder lifting otherwise-subthreshold arms into a functional regime and the switch then letting the same pre-positioned hub express coincidence suppression, \emph{provided} that suppression rests on the synchronous recruitment of a correlated hub population.

\begin{figure}[h]
\centering
\includegraphics[width=\linewidth]{robustness_sweep_en.png}
\caption{\textbf{XOR requires a correlated hub (Brian2, 100 seeds per point).} As the fraction of \emph{shared} afferent drive to the 20 hub cells falls from 1.0 (all hubs see the same two Poisson trains) towards 0.5 (each hub also receives its own independent, pattern-driven train), coincidence suppression degrades smoothly. \textbf{(A)} the fraction of runs classified XOR falls from $1.00$ through a knee near $\mathrm{shared}\!\approx\!0.82$, while the fraction in which \emph{both} direct arms grew stays at $1.00$ throughout. \textbf{(B)} the continuous selectivity index $\mathrm{xor\_index}=(\min(R_{10},R_{01})-R_{11})/\min(R_{10},R_{01})$ (median, IQR band) declines monotonically and crosses zero (XOR$\to$OR) near $\mathrm{shared}\!\approx\!0.66$. \textbf{(C)} the coincidence read-out rate $R_{11}$ rises from ${\sim}8$ to ${\sim}87$~Hz. Open triangles: a fully-random (unbalanced) pattern-exposure control at $\mathrm{shared}=1.0$, plotted at the value it took under the former absolute-rate classification (XOR $0.96$); the four exceptions are runs that fell below that $25$~Hz gate, and their continuous index was XOR-like throughout (\S\ref{app:growth}). On either classification the result does not depend on balanced exposure.}
\label{fig:robust}
\end{figure}

\FloatBarrier
One methodological point emerged from the set-point sweep and applies to every classified figure above. Our original classifier gated on an absolute read-out rate, calling a run \emph{collapse} when $\min(R_{10},R_{01})<25$~Hz; no figure in this preprint uses that criterion. That gate conflates two different failures: a builder that never taught, and a taught circuit that simply operates at a lower rate. All twenty successful runs at $\ECl^{\mathrm{build}}=-45$~mV were initially labelled collapse despite arms well above the cliff and $\mathrm{xor\_index}$ reaching $1.00$, purely because their single-input rates sat between $15$ and $24$~Hz. We have re-scored every run on \emph{absolute-rate-threshold-free} criteria, and the qualifier is deliberate: the scheme still uses thresholds, but circuit-intrinsic ones rather than a firing rate in hertz. Explicitly, with $R_{ij}$ the read-out rate for input pattern $ij$ and
\[
\mathrm{xor\_index}=\frac{\min(R_{10},R_{01})-R_{11}}{\min(R_{10},R_{01})},
\]
a run is labelled as follows. \emph{Failure to teach} when neither arm has cleared the firing cliff, $\max(w_I,w_P)\leq0.21$; \emph{asymmetric} when exactly one has, a case that with only two arms a mean or median would hide, so that the gate is applied to the individual weights and not to their average; \emph{silent} when both have cleared it but the read-out does not respond to either single input, $\min(R_{10},R_{01})=0$; and otherwise \emph{XOR} when $\mathrm{xor\_index}\geq0.5$ and \emph{OR} when it does not. The index is evaluated only in that last case, since $\min(R_{10},R_{01})>0$ is required for its denominator. The cliff is a property of the circuit, measurable independently of any run, and the index is dimensionless. Read-out rates are reported as descriptors and never enter the classification. Applying the cliff to the individual arms rather than to their mean is the correct criterion but changes little here: two runs of the $800$ are reclassified as \emph{asymmetric}, both in cells whose median weight sits barely above the cliff. Two further cells, at $(t_{\mathrm{build}},w_{\mathrm{HL}})=(50,16)$ and $(65,34)$, split their twenty seeds exactly ten to ten between XOR and OR; we mark them as ties rather than assign a majority, since a tie is a result about the cell and not a detail to resolve by convention. The re-scoring script is in the repository so that the two labellings can be compared directly.

Honest limits: the emergent XOR is clean and symmetric but \emph{soft} (coincidence held to a fraction of the single-input rate, not annihilated); the grown circuit sits in a low-rate regime; the plasticity rule, though BCM-grounded, is a stylized abstraction, not a fit to measured plasticity; both chloride set points are now mapped, the immature one bounded from below by the spike threshold and the mature one robust across the range tested (\S\ref{app:switchsim}); and the hub population size, twenty cells, is unconstrained by data. Above all, the architecture is preconfigured, so that what the model establishes is the necessity of the developmental \emph{sequence} for reaching a functional operating point, and not the emergence of the topology itself. Both simulations and every figure script are in the code repository (Data and code, main text).


\FloatBarrier
\section{The GABA switch in the hippocampus}\label{app:hipp}
In the mammalian hippocampus the switch is elaborated into a developmental \emph{programme}, and it is here that its constructive role is clearest. Inhibition is built first: GABAergic neurons and their synapses appear before glutamatergic ones, and in the immature network GABA is depolarizing and, under the right conditions, excitatory (Ben-Ari, 2002; Ben-Ari et al., 2007). The molecular cause is the late induction of the neuron-specific chloride extruder KCC2, which lowers $[\Clm]_{\mathrm{in}}$ and renders GABA hyperpolarizing as neurons mature; suppressing KCC2 shifts the GABA reversal potential back towards depolarizing values (Rivera et al., 1999); the family and its developmental regulation are reviewed by Kaila et al.\ (2014). Throughout this appendix we present the constructive reading of the immature depolarizing phase; the in vivo evidence for and against it is weighed in \S\ref{sec:switch}, and a reader who rejects that reading should take this appendix as describing a phenomenon of the slice.

While it remains depolarizing, GABA does not merely coexist with development: it \emph{drives} it. Immature GABA generates Giant Depolarizing Potentials (GDPs), the first coordinated network activity of the hippocampus (Ben-Ari et al., 1989), produced by the synergistic excitatory action of GABA$_{\mathrm{A}}$ and NMDA receptors and accompanied by the calcium transients that guide the growth and stabilization of synapses (Leinekugel et al., 1997). These events are orchestrated by a small population of richly connected GABAergic \emph{hub} interneurons, whose wide axonal fields pace the synchrony; stimulating a single hub can advance, delay, or abort a GDP across the whole network (Bonifazi et al., 2009; Cossart \& Khazipov, 2022). The picture is that of a common inhibitory hub which, in its depolarizing phase, coactivates and so helps to \emph{grow} the very excitatory arms it will later regulate.

The decisive subtlety is that the inhibitory \emph{conductance} is present the whole time; only the \emph{sign of its driving force} changes. As KCC2 rises and $\ECl$ falls towards and below rest, the same hub stops depolarizing and becomes inhibitory. Near rest the effect is predominantly shunting and divisive; farther below rest a hyperpolarizing component is added, as described in the biophysical background (Appendix~\ref{app:sim}). One topology; two developmental modes, selected by a slow chloride set point (Fig.~\ref{fig:dev}; the circuit itself, in its builder and governor forms, is Fig.~\ref{fig:devseq}A,C). And because KCC2 expression and transport activity can themselves be activity-regulated, the switch provides a plausible homeostatic negative-feedback route out of the highly synchronous GDP regime and towards the sparse, marginally stable operating point of the adult, the regime in which cortical circuits display scale-free avalanche statistics (Beggs \& Plenz, 2003) and towards which they are actively and homeostatically tuned in vivo (Ma et al., 2019).

\begin{figure}[h]
\centering
\includegraphics[width=0.92\linewidth]{preprint_fig_en.png}
\caption{\textbf{The builder$\rightarrow$governor hypothesis.} Early in the development the GABA driving force is strongly depolarizing (red); the inhibitory hub acts cooperatively and paces the giant depolarizing potentials (blue), whose calcium transients grow the glutamatergic arms of the circuit. As KCC2 rises and intracellular chloride falls (the ``GABA switch''), the driving force moves towards the mature inhibitory range, the GDPs disappear, and the same hub becomes an inhibitory governor into which the mature signed-XOR read-out is folded. Near rest that governor is shunting; below rest it also hyperpolarizes. One topology; two developmental modes, selected by a slow chloride set point. Schematic; axes illustrative.}
\label{fig:dev}
\end{figure}

\FloatBarrier
\section{The GABA switch in \emph{C.~elegans}}\label{app:worm}
The developmental GABA switch, once considered a vertebrate peculiarity, reaches down to the nematode. In the locomotor circuit of \emph{C.~elegans}, the GABA$_{\mathrm{A}}$ agonist muscimol \emph{excites} muscle in newly hatched (L1) animals and, by the middle of the first larval stage (about six hours after hatching), \emph{relaxes} it, a reversal from depolarizing to hyperpolarizing. The switch runs on the same chloride machinery as in mammals: the importer NKCC-1 sustains the early depolarizing response, while the extruders KCC-2 and ABTS-1 produce the later hyperpolarizing one. The switch still occurs (with a delay) in \emph{unc-25} mutants that cannot synthesize GABA at all, showing that in the worm it is largely a \emph{cell-intrinsic, programmed} maturation of the chloride gradient rather than something driven by neural activity (Han et al., 2015). The extruders were characterized independently: KCC-2 sets the chloride gradients required for inhibition, is upregulated during synaptic development, and operates in both neurons and muscle (Tanis et al., 2009); it acts together with ABTS-1, and when both are removed the gradient reverses so that GABA excites rather than inhibits (Bellemer et al., 2011).

A very recent result carries GABAergic signalling, and the chloride logic that sets its sign, down to the embryo. Immature GABAergic DD motor neurons transiently \emph{inhibit} early embryonic movement, about 8.5--9.5 hours after fertilization, \emph{before} their neurites have finished growing and before most of the nerve-ring synapses exist; the effect requires GABA synthesis in the DD neurons and UNC-49 GABA$_{\mathrm{A}}$ receptors in the muscle, occurs through both synaptic and non-synaptic (tonic) release, and follows chloride logic: removing the importer NKCC-1 deepens the inhibition while removing the extruder KCC-2 relieves it (Marvel-Coen et al., 2026). Two qualifications matter for our purposes, and we state them because it would be easy to over-read the result. The phenotype is revealed in \emph{snf-11} (GAT1) mutants, which lack the plasma-membrane GABA re-uptake transporter and therefore carry exaggerated GABA signalling, rather than in the wild type. And the action reported is \emph{inhibitory}, so the result does not extend the depolarizing builder phase into the worm embryo; if anything it indicates that GABA is already hyperpolarizing in embryonic muscle at that stage. What it does establish, and what we take from it, is that the transporter-controlled sign of GABA is a live and manipulable variable from the embryo onwards.

Two complications refine, rather than break, the picture. The worm has a \emph{second} route to an excitatory GABA that bypasses the chloride gradient entirely: EXP-1 is a GABA-gated cation channel, excitatory by construction (Beg \& Jorgensen, 2003), showing that the sign of a ``GABA'' signal can be set either by the transporter-controlled driving force or intrinsically by the receptor. And although eight connectomes have already been reconstructed across postembryonic development, revealing a largely constant wiring scaffold that becomes progressively more feedforward with age (Witvliet et al., 2021), connectomes report anatomy, not $\ECl$, and the ``builder'' depolarizing phase is largely embryonic. The worm establishes that the switch is conserved and cell-intrinsic; it does not, on its own, show the switch operating inside a central computational circuit. That is where the mammalian data (Appendix~\ref{app:hipp}) go further.

\FloatBarrier
\section{Chloride and GABA in plants}\label{app:plants}
A caveat, on which the main text also insists, comes before drawing the parallel: what follows is a claim of \emph{functional repurposing}, not of homology: the plant and animal proteins are unrelated, and all that is shared is the ancient \emph{principle} of ion control and the cotransporter family that implements it.

Long before neurons, cells used chloride to move. Chloride is the most abundant inorganic anion in plant cells and, though long dismissed as a mere micronutrient, it accumulates to osmotically decisive concentrations and is today recognized as a beneficial macronutrient, central to osmoregulation and turgor (Geilfus, 2018; Colmenero-Flores et al., 2019). The direction is what matters: plant cells keep $[\Clm]_{\mathrm{in}}$ \emph{high}, because in a walled cell a high internal osmolyte is what generates the turgor pressure that drives growth and movement, exactly the opposite of the mature neuron, which spends energy to keep $[\Clm]_{\mathrm{in}}$ \emph{low}.

Because turgor is the plant's actuator, chloride is the plant currency of movement. The opening and closing of stomata is a turgor motor in which anion channels are the active elements: the malate-activated vacuolar channel AtALMT9 is required for opening, and the R-type channel AtALMT12/QUAC1 mediates the anionic efflux of closure (De Angeli et al., 2013; Meyer et al., 2010). Faster movements use the same anions as a genuine excitable system: the Venus flytrap fires an all-or-none action potential carried by a calcium-activated \emph{anionic} conductance rather than a sodium current, and even \emph{counts} those action potentials to commit to digestion (B\"ohm et al., 2016). Plants even use GABA as a signal, but through a different protein and with the sign reversed: GABA regulates the ALMT anion channels, which are \emph{not} homologous to the animal GABA$_{\mathrm{A}}$ receptor (they share only a $\sim$12-amino-acid motif; Ramesh et al., 2015, 2017, 2018), and because the internal anions are high, opening an ALMT channel \emph{depolarizes} while GABA, by \emph{closing} it, hyperpolarizes.

Do plants have the switch itself? They have the machinery family but not the neuronal specialization. Cation-chloride cotransporters exist in plants, but flowering plants typically carry a single such gene (\emph{AtCCC} in Arabidopsis), phylogenetically closest to the KCC clade and yet localized largely to endomembranes and the vasculature, and not deployed as a plasma-membrane extruder that sets a resting $\ECl$ for inhibition (Colmenero-Flores et al., 2007; Henderson et al., 2018). Plants therefore possess every ancient ingredient: chloride, anion channels, GABA, cotransporters, but not the diversified extrusion switch, and their chloride ``default'' points in the wrong direction for inhibition. That missing switch is exactly what the animal lineage adds, and it is where movement becomes control.

\FloatBarrier
\section*{References}
\small
\begin{apabib}
\item Akerman, C.~J., \& Cline, H.~T. (2006). Depolarizing GABAergic conductances regulate the balance of excitation to inhibition in the developing retinotectal circuit in vivo. \emph{Journal of Neuroscience, 26}(19), 5117--5130. \url{https://doi.org/10.1523/JNEUROSCI.0319-06.2006}
\item Beg, A.~A., \& Jorgensen, E.~M. (2003). EXP-1 is an excitatory GABA-gated cation channel. \emph{Nature Neuroscience, 6}(11), 1145--1152. \url{https://doi.org/10.1038/nn1136}
\item Beggs, J.~M., \& Plenz, D. (2003). Neuronal avalanches in neocortical circuits. \emph{Journal of Neuroscience, 23}(35), 11167--11177. \url{https://doi.org/10.1523/JNEUROSCI.23-35-11167.2003}
\item Bellemer, A., Hirata, T., Romero, M.~F., \& Koelle, M.~R. (2011). Two types of chloride transporters are required for GABA$_{\mathrm{A}}$ receptor-mediated inhibition in \emph{C.~elegans}. \emph{The EMBO Journal, 30}(9), 1852--1863. \url{https://doi.org/10.1038/emboj.2011.83}
\item Ben-Ari, Y. (2002). Excitatory actions of GABA during development: the nature of the nurture. \emph{Nature Reviews Neuroscience, 3}(9), 728--739. \url{https://doi.org/10.1038/nrn920}
\item Ben-Ari, Y. (2014). The GABA excitatory/inhibitory developmental sequence: a personal journey. \emph{Neuroscience, 279}, 187--219. \url{https://doi.org/10.1016/j.neuroscience.2014.08.001}
\item Ben-Ari, Y., Cherubini, E., Corradetti, R., \& Gaiarsa, J.-L. (1989). Giant synaptic potentials in immature rat CA3 hippocampal neurones. \emph{The Journal of Physiology, 416}, 303--325. \url{https://doi.org/10.1113/jphysiol.1989.sp017762}
\item Ben-Ari, Y., Gaiarsa, J.-L., Tyzio, R., \& Khazipov, R. (2007). GABA: a pioneer transmitter that excites immature neurons and generates primitive oscillations. \emph{Physiological Reviews, 87}(4), 1215--1284. \url{https://doi.org/10.1152/physrev.00017.2006}
\item Bienenstock, E.~L., Cooper, L.~N., \& Munro, P.~W. (1982). Theory for the development of neuron selectivity: orientation specificity and binocular interaction in visual cortex. \emph{Journal of Neuroscience, 2}(1), 32--48. \url{https://doi.org/10.1523/JNEUROSCI.02-01-00032.1982}
\item B\"ohm, J., Scherzer, S., Krol, E., Kreuzer, I., von Meyer, K., Lorey, C., \ldots\ Hedrich, R. (2016). The Venus flytrap \emph{Dionaea muscipula} counts prey-induced action potentials to induce sodium uptake. \emph{Current Biology, 26}(3), 286--295. \url{https://doi.org/10.1016/j.cub.2015.11.057}
\item Bonifazi, P., Goldin, M., Picardo, M.~A., Jorquera, I., Cattani, A., Bianconi, G., Represa, A., Ben-Ari, Y., \& Cossart, R. (2009). GABAergic hub neurons orchestrate synchrony in developing hippocampal networks. \emph{Science, 326}(5958), 1419--1424. \url{https://doi.org/10.1126/science.1175509}
\item Bregestovski, P., \& Bernard, C. (2012). Excitatory GABA: how a correct observation may turn out to be an experimental artifact. \emph{Frontiers in Pharmacology, 3}, 65. \url{https://doi.org/10.3389/fphar.2012.00065}
\item Cancedda, L., Fiumelli, H., Chen, K., \& Poo, M.-m. (2007). Excitatory GABA action is essential for morphological maturation of cortical neurons in vivo. \emph{Journal of Neuroscience, 27}(19), 5224--5235. \url{https://doi.org/10.1523/JNEUROSCI.5169-06.2007}
\item Carandini, M., \& Heeger, D.~J. (2012). Normalization as a canonical neural computation. \emph{Nature Reviews Neuroscience, 13}(1), 51--62. \url{https://doi.org/10.1038/nrn3136}
\item Caz\'e, R.~D., Humphries, M., \& Gutkin, B. (2013). Passive dendrites enable single neurons to compute linearly non-separable functions. \emph{PLoS Computational Biology, 9}(2), e1002867. \url{https://doi.org/10.1371/journal.pcbi.1002867}
\item Chance, F.~S., Abbott, L.~F., \& Reyes, A.~D. (2002). Gain modulation from background synaptic input. \emph{Neuron, 35}(4), 773--782. \url{https://doi.org/10.1016/S0896-6273(02)00820-6}
\item Clopath, C., B\"using, L., Vasilaki, E., \& Gerstner, W. (2010). Connectivity reflects coding: a model of voltage-based STDP with homeostasis. \emph{Nature Neuroscience, 13}(3), 344--352. \url{https://doi.org/10.1038/nn.2479}
\item Colmenero-Flores, J.~M., Mart\'inez, G., Gamba, G., V\'azquez, N., Iglesias, D.~J., Brum\'os, J., \& Tal\'on, M. (2007). Identification and functional characterization of cation--chloride cotransporters in plants. \emph{The Plant Journal, 50}(2), 278--292. \url{https://doi.org/10.1111/j.1365-313X.2007.03048.x}
\item Colmenero-Flores, J.~M., Franco-Navarro, J.~D., Cubero-Font, P., Peinado-Torrubia, P., \& Rosales, M.~A. (2019). Chloride as a beneficial macronutrient in higher plants. \emph{International Journal of Molecular Sciences, 20}(19), 4686. \url{https://doi.org/10.3390/ijms20194686}
\item Cossart, R., \& Khazipov, R. (2022). How development sculpts hippocampal circuits and function. \emph{Physiological Reviews, 102}(1), 343--378. \url{https://doi.org/10.1152/physrev.00044.2020}
\item De Angeli, A., Zhang, J., Meyer, S., \& Martinoia, E. (2013). AtALMT9 is a malate-activated vacuolar chloride channel required for stomatal opening in Arabidopsis. \emph{Nature Communications, 4}, 1804. \url{https://doi.org/10.1038/ncomms2815}
\item Durand, G.~M., Kovalchuk, Y., \& Konnerth, A. (1996). Long-term potentiation and functional synapse induction in developing hippocampus. \emph{Nature, 381}(6577), 71--75. \url{https://doi.org/10.1038/381071a0}
\item Dzhala, V.~I., Talos, D.~M., Sdrulla, D.~A., Brumback, A.~C., Mathews, G.~C., Benke, T.~A., Delpire, E., Jensen, F.~E., \& Staley, K.~J. (2005). NKCC1 transporter facilitates seizures in the developing brain. \emph{Nature Medicine, 11}(11), 1205--1213. \url{https://doi.org/10.1038/nm1301}
\item Friedel, P., Kahle, K.~T., Zhang, J., Hertz, N., Pisella, L.~I., Buhler, E., Schaller, F., Duan, J., Khanna, A.~R., Bishop, P.~N., Shokat, K.~M., \& Medina, I. (2015). WNK1-regulated inhibitory phosphorylation of the KCC2 cotransporter maintains the depolarizing action of GABA in immature neurons. \emph{Science Signaling, 8}(383), ra65. \url{https://doi.org/10.1126/scisignal.aaa0354}
\item Geilfus, C.-M. (2018). Chloride: from nutrient to toxicant. \emph{Plant and Cell Physiology, 59}(5), 877--886. \url{https://doi.org/10.1093/pcp/pcy071}
\item Gidon, A., Zolnik, T.~A., Fidzinski, P., Bolduan, F., Papoutsi, A., Poirazi, P., Holtkamp, M., Vida, I., \& Larkum, M.~E. (2020). Dendritic action potentials and computation in human layer 2/3 cortical neurons. \emph{Science, 367}(6473), 83--87. \url{https://doi.org/10.1126/science.aax6239}
\item Graupner, M., \& Brunel, N. (2012). Calcium-based plasticity model explains sensitivity of synaptic changes to spike pattern, rate, and dendritic location. \emph{Proceedings of the National Academy of Sciences, 109}(10), 3991--3996. \url{https://doi.org/10.1073/pnas.1109359109}
\item Han, B., Bellemer, A., \& Koelle, M.~R. (2015). An evolutionarily conserved switch in response to GABA affects development and behavior of the locomotor circuit of \emph{Caenorhabditis elegans}. \emph{Genetics, 199}(4), 1159--1172. \url{https://doi.org/10.1534/genetics.114.173963}
\item Hartmann, A.-M., Tesch, D., Nothwang, H.~G., \& Bininda-Emonds, O.~R.~P. (2014). Evolution of the cation chloride cotransporter family: ancient origins, gene losses, and subfunctionalization through duplication. \emph{Molecular Biology and Evolution, 31}(2), 434--447. \url{https://doi.org/10.1093/molbev/mst225}
\item Henderson, S.~W., Wege, S., \& Gilliham, M. (2018). Plant cation-chloride cotransporters (CCC): evolutionary origins and functional insights. \emph{International Journal of Molecular Sciences, 19}(2), 492. \url{https://doi.org/10.3390/ijms19020492}
\item Holt, G.~R., \& Koch, C. (1997). Shunting inhibition does not have a divisive effect on firing rates. \emph{Neural Computation, 9}(5), 1001--1013. \url{https://doi.org/10.1162/neco.1997.9.5.1001}
\item Isaac, J.~T.~R., Nicoll, R.~A., \& Malenka, R.~C. (1995). Evidence for silent synapses: implications for the expression of LTP. \emph{Neuron, 15}(2), 427--434. \url{https://doi.org/10.1016/0896-6273(95)90046-2}
\item J\'ekely, G., Keijzer, F., \& Godfrey-Smith, P. (2015). An option space for early neural evolution. \emph{Philosophical Transactions of the Royal Society B, 370}(1684), 20150181. \url{https://doi.org/10.1098/rstb.2015.0181}
\item Kaila, K. (1994). Ionic basis of GABA$_{\mathrm{A}}$ receptor channel function in the nervous system. \emph{Progress in Neurobiology, 42}(4), 489--537. \url{https://doi.org/10.1016/0301-0082(94)90049-3}
\item Kaila, K., Price, T.~J., Payne, J.~A., Puskarjov, M., \& Voipio, J. (2014). Cation-chloride cotransporters in neuronal development, plasticity and disease. \emph{Nature Reviews Neuroscience, 15}(10), 637--654. \url{https://doi.org/10.1038/nrn3819}
\item Keijzer, F., van Duijn, M., \& Lyon, P. (2013). What nervous systems do: early evolution, input--output, and the skin brain thesis. \emph{Adaptive Behavior, 21}(2), 67--85. \url{https://doi.org/10.1177/1059712312465330}
\item Kirmse, K., Kummer, M., Kovalchuk, Y., Witte, O.~W., Garaschuk, O., \& Holthoff, K. (2015). GABA depolarizes immature neurons and inhibits network activity in the neonatal neocortex in vivo. \emph{Nature Communications, 6}, 7750. \url{https://doi.org/10.1038/ncomms8750}
\item Leinekugel, X., Medina, I., Khalilov, I., Ben-Ari, Y., \& Khazipov, R. (1997). Ca$^{2+}$ oscillations mediated by the synergistic excitatory actions of GABA$_{\mathrm{A}}$ and NMDA receptors in the neonatal hippocampus. \emph{Neuron, 18}(2), 243--255. \url{https://doi.org/10.1016/S0896-6273(00)80265-2}
\item Liao, D., Hessler, N.~A., \& Malinow, R. (1995). Activation of postsynaptically silent synapses during pairing-induced LTP in CA1 region of hippocampal slice. \emph{Nature, 375}(6530), 400--404. \url{https://doi.org/10.1038/375400a0}
\item Ma, Z., Turrigiano, G.~G., Wessel, R., \& Hengen, K.~B. (2019). Cortical circuit dynamics are homeostatically tuned to criticality in vivo. \emph{Neuron, 104}(4), 655--664.e4. \url{https://doi.org/10.1016/j.neuron.2019.08.031}
\item Marco de Lucas, J. (2024). From worms to mice: homeostasis maybe all you need. \emph{arXiv}. \url{https://arxiv.org/abs/2412.20090}
\item Marco de Lucas, J., Pe\~na Fern\'andez, M., \& Lloret Iglesias, L. (2024). Implementing engrams from a machine learning perspective: XOR as a basic motif. \emph{arXiv}. \url{https://arxiv.org/abs/2406.09940}
\item Marco de Lucas, J., Pe\~na Fern\'andez, M., \& Lloret Iglesias, L. (2026). \emph{Signed XOR feedback} (European Patent Application No.~26382252.0). European Patent Office.
\item Marvel-Coen, J., Ardiel, E., Zhao, J., Nurrish, S., \& Kaplan, J.~M. (2026). Immature \emph{Caenorhabditis elegans} motor neurons control early embryo behavior via both synaptic and nonsynaptic GABA release. \emph{Proceedings of the National Academy of Sciences, 123}(2), e2508541123. \url{https://doi.org/10.1073/pnas.2508541123}
\item Meyer, S., Mumm, P., Imes, D., Endler, A., Weder, B., Al-Rasheid, K.~A.~S., \ldots\ Hedrich, R. (2010). AtALMT12 represents an R-type anion channel required for stomatal movement in Arabidopsis guard cells. \emph{The Plant Journal, 63}(6), 1054--1062. \url{https://doi.org/10.1111/j.1365-313X.2010.04302.x}
\item Moroz, L.~L., Kocot, K.~M., Citarella, M.~R., Dosung, S., Norekian, T.~P., Povolotskaya, I.~S., \ldots\ Kohn, A.~B. (2014). The ctenophore genome and the evolutionary origins of neural systems. \emph{Nature, 510}(7503), 109--114. \url{https://doi.org/10.1038/nature13400}
\item Murata, Y., \& Colonnese, M.~T. (2020). GABAergic interneurons excite neonatal hippocampus in vivo. \emph{Science Advances, 6}(24), eaba1430. \url{https://doi.org/10.1126/sciadv.aba1430}
\item Pe\~na Fern\'andez, M., Gonz\'alez R\'ios, A., Lloret Iglesias, L., \& Marco de Lucas, J. (2026a). From homeostasis to credit assignment: a signed-XOR connectomic motif for local directional error signalling. \emph{bioRxiv}. \url{https://doi.org/10.64898/2026.06.05.730322}
\item Pe\~na Fern\'andez, M., Lloret Iglesias, L., \& Marco de Lucas, J. (2026b). Toward defining loss functions in neuroscience: an XOR-based neuronal mechanism. \emph{bioRxiv}. \url{https://doi.org/10.64898/2026.03.16.712061}
\item Pe\~na Fern\'andez, M., Lloret Iglesias, L., \& Marco de Lucas, J. (2026c). Signed-XOR error and sparse coding in a Dale-compliant substrate for sequence memorization. \emph{bioRxiv}. \url{https://doi.org/10.64898/2026.06.24.734176}
\item Piala, A.~T., Moon, T.~M., Akella, R., He, H., Cobb, M.~H., \& Goldsmith, E.~J. (2014). Chloride sensing by WNK1 involves inhibition of autophosphorylation. \emph{Science Signaling, 7}(324), ra41. \url{https://doi.org/10.1126/scisignal.2005050}
\item Pouille, F., \& Scanziani, M. (2001). Enforcement of temporal fidelity in pyramidal cells by somatic feed-forward inhibition. \emph{Science, 293}(5532), 1159--1163. \url{https://doi.org/10.1126/science.1060342}
\item Prescott, S.~A., \& De Koninck, Y. (2003). Gain control of firing rate by shunting inhibition: roles of synaptic noise and dendritic saturation. \emph{Proceedings of the National Academy of Sciences, 100}(4), 2076--2081. \url{https://doi.org/10.1073/pnas.0337591100}
\item Ramesh, S.~A., Tyerman, S.~D., Xu, B., Bose, J., Kaur, S., Conn, V., \ldots\ Gilliham, M. (2015). GABA signalling modulates plant growth by directly regulating the activity of plant-specific anion transporters. \emph{Nature Communications, 6}, 7879. \url{https://doi.org/10.1038/ncomms8879}
\item Ramesh, S.~A., Tyerman, S.~D., Gilliham, M., \& Xu, B. (2017). $\gamma$-Aminobutyric acid (GABA) signalling in plants. \emph{Cellular and Molecular Life Sciences, 74}(9), 1577--1603. \url{https://doi.org/10.1007/s00018-016-2415-7}
\item Ramesh, S.~A., Kamran, M., Sullivan, W., Chirkova, L., Okamoto, M., Degryse, F., McLaughlin, M., Gilliham, M., \& Tyerman, S.~D. (2018). Aluminum-activated malate transporters can facilitate GABA transport. \emph{The Plant Cell, 30}(5), 1147--1164. \url{https://doi.org/10.1105/tpc.17.00864}
\item Rheims, S., Holmgren, C.~D., Chazal, G., Mulder, J., Harkany, T., Zilberter, T., \& Zilberter, Y. (2009). GABA action in immature neocortical neurons directly depends on the availability of ketone bodies. \emph{Journal of Neurochemistry, 110}(4), 1330--1338. \url{https://doi.org/10.1111/j.1471-4159.2009.06230.x}
\item Rinehart, J., Maksimova, Y.~D., Tanis, J.~E., Stone, K.~L., Hodson, C.~A., Zhang, J., Risinger, M., Pan, W., Wu, D., Colangelo, C.~M., Forbush, B., Joiner, C.~H., Gulcicek, E.~E., Gallagher, P.~G., \& Lifton, R.~P. (2009). Sites of regulated phosphorylation that control K-Cl cotransporter activity. \emph{Cell, 138}(3), 525--536. \url{https://doi.org/10.1016/j.cell.2009.05.031}
\item Rivera, C., Voipio, J., Payne, J.~A., Ruusuvuori, E., Lahtinen, H., Lamsa, K., Pirvola, U., Saarma, M., \& Kaila, K. (1999). The K$^{+}$/Cl$^{-}$ co-transporter KCC2 renders GABA hyperpolarizing during neuronal maturation. \emph{Nature, 397}(6716), 251--255. \url{https://doi.org/10.1038/16697}
\item Ruusuvuori, E., Kirilkin, I., Pandya, N., \& Kaila, K. (2010). Spontaneous network events driven by depolarizing GABA action in neonatal hippocampal slices are not attributable to deficient mitochondrial energy metabolism. \emph{Journal of Neuroscience, 30}(46), 15638--15642. \url{https://doi.org/10.1523/JNEUROSCI.3355-10.2010}
\item Scellier, B., \& Bengio, Y. (2017). Equilibrium propagation: bridging the gap between energy-based models and backpropagation. \emph{Frontiers in Computational Neuroscience, 11}, 24. \url{https://doi.org/10.3389/fncom.2017.00024}
\item Shatz, C.~J. (1992). The developing brain. \emph{Scientific American, 267}(3), 60--67. \url{https://doi.org/10.1038/scientificamerican0992-60}
\item Silver, R.~A. (2010). Neuronal arithmetic. \emph{Nature Reviews Neuroscience, 11}(7), 474--489. \url{https://doi.org/10.1038/nrn2864}
\item Tanis, J.~E., Bellemer, A., Moresco, J.~J., Forbush, B., \& Koelle, M.~R. (2009). The potassium chloride cotransporter KCC-2 coordinates development of inhibitory neurotransmission and synapse structure in \emph{Caenorhabditis elegans}. \emph{Journal of Neuroscience, 29}(32), 9943--9954. \url{https://doi.org/10.1523/JNEUROSCI.1989-09.2009}
\item Turrigiano, G.~G. (2008). The self-tuning neuron: synaptic scaling of excitatory synapses. \emph{Cell, 135}(3), 422--435. \url{https://doi.org/10.1016/j.cell.2008.10.008}
\item Tyzio, R., Cossart, R., Khalilov, I., Minlebaev, M., H\"ubner, C.~A., Represa, A., Ben-Ari, Y., \& Khazipov, R. (2006). Maternal oxytocin triggers a transient inhibitory switch in GABA signaling in the fetal brain during delivery. \emph{Science, 314}(5806), 1788--1792. \url{https://doi.org/10.1126/science.1133212}
\item Tzilivaki, A., Kastellakis, G., \& Poirazi, P. (2019). Challenging the point neuron dogma: FS basket cells as 2-stage nonlinear integrators. \emph{Nature Communications, 10}, 3664. \url{https://doi.org/10.1038/s41467-019-11537-7}
\item Valeeva, G., Tressard, T., Mukhtarov, M., Baude, A., \& Khazipov, R. (2016). An optogenetic approach for investigation of excitatory and inhibitory network GABA actions in mice expressing channelrhodopsin-2 in GABAergic neurons. \emph{Journal of Neuroscience, 36}(22), 5961--5973. \url{https://doi.org/10.1523/JNEUROSCI.3482-15.2016}
\item Vogels, T.~P., Sprekeler, H., Zenke, F., Clopath, C., \& Gerstner, W. (2011). Inhibitory plasticity balances excitation and inhibition in sensory pathways and memory networks. \emph{Science, 334}(6062), 1569--1573. \url{https://doi.org/10.1126/science.1211095}
\item Wang, D.~D., \& Kriegstein, A.~R. (2008). GABA regulates excitatory synapse formation in the neocortex via NMDA receptor activation. \emph{Journal of Neuroscience, 28}(21), 5547--5558. \url{https://doi.org/10.1523/JNEUROSCI.5599-07.2008}
\item Whittington, J.~C.~R., \& Bogacz, R. (2017). An approximation of the error backpropagation algorithm in a predictive coding network with local Hebbian synaptic plasticity. \emph{Neural Computation, 29}(5), 1229--1262. \url{https://doi.org/10.1162/NECO\_a\_00949}
\item Witvliet, D., Mulcahy, B., Mitchell, J.~K., Meirovitch, Y., Berger, D.~R., Wu, Y., \ldots\ Zhen, M. (2021). Connectomes across development reveal principles of brain maturation. \emph{Nature, 596}(7871), 257--261. \url{https://doi.org/10.1038/s41586-021-03778-8}
\item Zhang, J., Cordshagen, A., Medina, I., Nothwang, H.~G., Wisniewski, J.~R., Winklhofer, M., \& Hartmann, A.-M. (2020). Staurosporine and NEM mainly impair WNK-SPAK/OSR1 mediated phosphorylation of KCC2 and NKCC1. \emph{PLoS ONE, 15}(5), e0232967. \url{https://doi.org/10.1371/journal.pone.0232967}
\item Zilberter, M. (2016). Reality of inhibitory GABA in neonatal brain: time to rewrite the textbooks? \emph{Journal of Neuroscience, 36}(40), 10242--10244. \url{https://doi.org/10.1523/JNEUROSCI.2270-16.2016}
\end{apabib}

\end{document}
